\documentclass{article}

\usepackage[utf8]{inputenc}
\usepackage[T1]{fontenc}
\usepackage{amsmath, amsthm}
\usepackage{amssymb}
\usepackage{stmaryrd}
\usepackage{listings}
\usepackage[svgnames]{xcolor}
\usepackage{tikz}
\usepackage{array}
\usepackage{graphicx}
\usepackage[backend=biber, maxbibnames=9]{biblatex}
\usepackage[noend]{algpseudocode}
\usepackage{algorithm}
\usepackage{mathtools}
\usepackage{subcaption}
\usepackage{wrapfig}
\usepackage{hhline}
\usepackage{bsymb}
\usepackage{longtable}
\usepackage{syntax}
\usepackage{ebproof}
\usepackage{minted}
\usepackage[left=1.5cm, right=1.5cm]{geometry}
\usepackage{pifont}% http://ctan.org/pkg/pifont
\usepackage{hyperref}

% NOTE:
% File has moved to overleaf for collaboration!!!

\usemintedstyle{tango} %or tango, vs, xcode, solarized-light
\setminted{bgcolor=gray!5, tabsize=4}
% Commands
% Simple set comprehension notation
\newcommand{\Set}[2]{%
  \{\, #1 \mid #2 \, \}%
}
% Event-B-like set notation
\newcommand{\bSet}[3]{%
  \{\, #1 \cdot #2 \mid #3 \, \}%
}
\newcommand{\bSetT}[2]{%
  \{\, #1 \cdot #2 \,\}%
}
\newcommand{\setOne}[1]{%
  \{\, #1 \,\}%
}

% Bag notation
% \newcommand{\lbbar}{\{\kern-0.5ex|}
% \newcommand{\rbbar}{|\kern-0.5ex\}}
\newcommand{\bag}[3]{%
  \llbracket \, #1 \cdot #2 \mid #3 \, \rrbracket%
}
\newcommand{\bagT}[2]{%
  \llbracket \, #1 \cdot #2 \,\rrbracket%
}
\newcommand{\Bag}[2]{%
  \llbracket \, #1 \mid #2 \, \rrbracket%
}
\newcommand{\bagOne}[1]{%
  \llbracket \, #1 \, \rrbracket%
}

% List comprehension notation
\newcommand{\List}[3]{%
  [\, #1 \cdot #2 \mid #3 \, ]%
}
\newcommand{\ListT}[2]{%
  [\, #1 \cdot #2 \, ]%
}
\newcommand{\ListG}[2]{%
  [\, #1 \mid #2 \, ]%
}
\newcommand{\listOne}[1]{%
  [\, #1 \, ]%
}
% List concatenation
\newcommand{\concat}{%
  \mathbin{{+}\mspace{-8mu}{+}}%
}
% \DeclareUnicodeCharacter{29FA}{\concat}

% Lex table entry
\newcommand{\lexentry}[3]{%
  \ensuremath{#1} & \texttt{#2} & #3\\%
}

\newcommand{\cmark}{\ding{51}}%
\newcommand{\xmark}{\ding{55}}%
\newcommand{\ccmark}{%
    \cmark\kern-0.4em\cmark%
}

% From https://tex.stackexchange.com/questions/82782/footnote-in-align-environment
\makeatletter
\let\original@footnote\footnote
\newcommand{\align@footnote}[1]{%
  \ifmeasuring@
    \chardef\@tempfn=\value{footnote}%
    \footnotemark
    \setcounter{footnote}{\@tempfn}%
  \else
    \iffirstchoice@
      \original@footnote{#1}%
    \fi
  \fi}
\pretocmd{\start@align}{\let\footnote\align@footnote}{}{}
\makeatother

% Modify appearance of nonterminals
% Enables italics for nonterminals
% \renewcommand{\syntleft}{\itshape}
\renewcommand{\syntleft}{}
\renewcommand{\syntright}{}

% Other preamble
\allowdisplaybreaks
\graphicspath{{./images/}}
% NOTE: need to run `biber <name>` to ensure up-to-date references
\addbibresource{../ref.bib}

\title{Complete Specification and Design for Simile}
\author{Anthony Hunt}

\begin{document}
\maketitle
\tableofcontents
\newpage

\section{High-Level Design}
\subsection{Overview}
This document serves as a living specification and design record of the very-high-level programming language Simile. Simile is founded on the union between formal specification languages and efficiency-focused programming languages, with the belief that a language can be both highly performant and highly abstract. We focus on providing a python-like/specification-language-like programming experience with a compiler strong enough to compete with hand-optimized C code. Where possible, we aim to formalize Simile's compilation process to ensure that aggressively optimized programs remain correct by construction.

% TODO dont talk about event b so much here - prime real estate to introduce Simile
Simile takes its syntactic and semantic roots in that of Event-B, a set-theory-first language used to formalize safety-oriented software specifications. Event-B notation often leans on set comprehensions (of the form $\bSet{x}{x \in S}{E}$) and set operations (like $S \cup T$), with progressive specification refinement at the core of the development process. Like object-oriented programs, Event-B programs tend to operate at multiple refinement levels simultaneously, especially when proving program correctness. Each atomic action must be proved correct, then proofs may be composed to prove larger programs, and so on. With Event-B as inspiration, Simile's syntax and semantics is intended to be proof-friendly in the same way, by following discrete mathematics notation.

At the same time, Simile attempts to use the semantic information available at high-level programs to aggressively optimize the executable value. For example, $X \cap T = T$ may require evaluation of $X \cap T$ in a normal programming language, but Simile attempts to match known theorems against expressions to determine that $X \cap T = T$ is only true when $X \supseteq T$. By avoiding the difficult reasoning paths of traditional imperative languages (involving side effects or mutable objects) to perform an action, set theory semantics affords an entire class of optimizations not available to other programming languages.

By building from 50 years of programming language and formal methods literature, we hope to create a language that will spark further interest in simple-to-use, efficient-to-run languages. Following sections include the high-level design of Simile, formalized subsystems, and an implementation reference.

\subsection{Scope}
In scope:
\begin{itemize}
    \item We follow set theory semantics most of the time, diverging only for a particularly useful expression or programming objective.
    \item Programs should remain as abstract as possible, but users may provide helpful information to the compiler to unlock more aggressive optimizations.
    \item Notation follows that of Event-B (mostly), including a rich set of set and relational operators.
\end{itemize}
%%
Out of scope:
\begin{itemize}
    \item We cannot hope to become the objectively best language in all use cases. Each language has their strength, and Simile aims to be well-rounded; ergonomic enough to pick up, and fast enough to be more usable than other specification language runtimes.
    \item We do not follow conventional imperative language semantics that would otherwise clash with set theory optimizations.
\end{itemize}

\subsection{Background}
\subsubsection{Programmers should not worry about implementation details}
In high level programming languages, users often reach for abstract constructs like relations (dictionaries in Python, maps in Java), lists, and sets, largely because they are standardized, concise, and expressive. For example, HashMaps can be found in almost every language's standard library, and it would be peculiar to find a language that does not have first-class support for list collections. Most of the time, programmers using these types don't think about the layout or execution of their data, just that it is held within a container and operable under several composable actions (concatenation, union, intersection, lookups, etc.). However, this abstraction under the semantics of traditional high-level imperative/OO programs can lead to wasted memory and computation time, especially when composing multiple operations together. For example, the expression $S \cap (A \cup B \cup C)$ is theoretically bounded by the size of $S$, but a conventional language would construct several intermediary sets for $A \cup B$, then $(A \cup B) \cup C$, consuming more memory and time than required.

Low-level languages share an adjacent problem: programmers may trust primitive data type operations to be fast, but they must also tune every piece of their code to fit the type's implementation. Formally simple and powerful operations like list concatenation force the user to consider implementation details like memory allocation and ownership for both inputs and outputs. For example, in $A := S \concat T$, does the specific implementation create a new list for $A$ containing a shallow copy of the contents of $S$ and $T$? Or perhaps $A$ and $S$ will reference the same object, where the data in $T$ is moved to the pointed-to value of $S$. Programmers must track every byte of data for straightforward operations, causing mental overhead, distractions from program correctness, and the risk of implementation errors.

Abstract data types are semantically powerful and have been implemented by hundreds of languages across thousands of libraries. Further, they can be represented by a variety of concrete types, with wildly varying levels of efficiency and effectiveness for different use cases. For example, a set of integers can be implemented by a HashSet, a bitset, an array, a tree set, or even a trie. HashSets are advantageous for element lookups, but may not be able to compare to a dense bitset in terms of memory efficiency. An array could be effective for a set that may need to be sorted, but doesn't have direct element lookups. With a little metadata-help from the programmer, the compiler should be intelligent enough to make decisions on the concrete type and optimize expressions tailored to the concrete type while completely following the abstract type's semantics.

\paragraph{Motivating Example: A Visitor Information System}
Let's turn towards an example for specification-focused programming: a group of conference administrators create a program to track resourcing for a week's worth of workshops. They already have a list of visitors, available rooms, and planned workshops for each room. A specification would model these records via carrier sets (i.e., enumerated sets):
\begin{gather*}
    Visitor = \{Ryan, Adrian, Kevin, Charlie, Michael, ...\}\\
    Room = \{100, 101, 201, ...\}\\
    Workshop = \{ESOP, FASE, FoSSaCS\}
\end{gather*}
To make this model more useful, we express relationships between these sets. For example, every workshop must occur in a room and visitors will attend workshops. Further, we can place restrictions on each relation concisely through specification language notation. If we know workshops occur in a room, it would be safe to tighten our observation that \textit{all} workshops must occur in one room \textit{at a time}. So, we can explicitly refine the relationship to an injective (one-to-one) function that is total on the domain (the carrier set $Workshop$). In formal notation, we write concisely:
\begin{gather*}
    location: Workshop \tinj Room
\end{gather*}
We can further refine another relationship: visitors can attend workshops, but visitors can only attend one workshop \textit{at a time} and \textit{not all} visitors may attend a workshop. This relation is actually partial on the carrier set of $Visitor$ and a functional (many-to-one) relation. Formally:
\begin{gather*}
    attends: Visitor \pfun Workshop
\end{gather*}
With these relations, we can query for interesting information concisely. Let's say we want to throw a surprise dinner for all attendees in a specific room. We can compose the $attends$ relation with the $location$ relation to find the number of people expected in a specific room. For a specific $room$, we write:
\begin{gather*}
    card((location^{-1} \circ attends^{-1})[\{room\}]
\end{gather*}
A HashMap would be the first concrete data structure of choice for most programmers, but without knowledge of the refined relationships, we cannot guarantee that the relational inverse operation is guaranteed to work. We set aside HashMaps for now for the sole reason that they are limited in their representation capabilities; we cannot model a many-to-many relationship using the convential definition of the type.

If we were to build a compiler for this specification language, a naive implementation might only see the broad connection between these sets and choose to safely model them as a HashSet of tuples (which can represent many-to-many relations). In this case, we concretize the types of $location$ and $attends$ like so:
\begin{gather*}
    location: HashSet[(Workshop, Room)]\\
    attends: HashSet[(Visitor, Workshop)]
\end{gather*}
Then, each operation from the above query can be ordered in an imperative fashion:
\begin{algorithmic}
    \State $locationInverse: HashSet[(Room, Workshop)] := Inverse(location)$
    \State $attendsInverse: HashSet[(Workshop, Visitor)] := Inverse(attends)$
    \State $visitorsInRooms: HashSet[(Room, Visitor)] := Compose(locationInverse, attendsInverse)$
    \State $visitorsInQueryRoom: HashSet[Visitor] := Lookup(room, visitorsInRooms)$
    \State $mealsToPrepForRoom: int := Card(visitorsInQueryRoom)$
\end{algorithmic}
With a library of concrete functions:
\begin{algorithmic}
    \Function{Inverse}{$r: HashSet[(X, Y)]$}
        \Comment{Runtime: $O(n)$. Memory use: $O(n)$}
        \State $r': HashSet[(Y, X)] := \{\}$
        \For{$x,y \in r$}
            \State $r'.add((y,x))$ \Comment{$s.add$ is a primitive HashSet operation mutating $s$}
        \EndFor
        \Return $r'$
    \EndFunction
    \Function{Compose}{$r_1: HashSet[(X, Y)], r_2: HashSet[(Y, Z)]$}
        \Comment{Runtime: $O(n^2)$. Memory use: $O(n)$}
        \State $r': HashSet[(X,Z)] := \{\}$
        \For{$x,y_1 \in r_1$}
            \For{$y_2,z \in r_2$}
                \If{$y_1 = y_2$}
                    \State $r'.add((x,y))$
                \EndIf
            \EndFor
        \EndFor
        \Return $r'$
    \EndFunction
    \Function{Lookup}{$searchFor: X, r: HashSet[(X,Y)]$}
        \Comment{Runtime: $O(n)$. Memory use: $O(n)$}
        \State $r': HashSet[Y] := \{\}$
        \For{$x,y \in r$}
            \If{$searchFor = x$}
                \State $r'.add(y)$
            \EndIf
        \EndFor
        \Return $r'$
    \EndFunction
    \Function{Card}{$s: HashSet$}
        \Comment{Runtime: $O(n)$. Memory use: $O(1)$}
        \State $c: int := 0$
        \For{$x \in s$}
            \State $c := c + 1$
        \EndFor
        \Return $c$
    \EndFunction
\end{algorithmic}
As evident by the poor runtime and memory complexity, this is far from the most efficient representation of the system. However, it is a direct, valid translation, fulfilling all semantic requirements. A human may never write such code, but a compiler without statically analyzing the use of abstract times, may choose the broadest representation at the cost of runtime and memory usage. With knowledge of the problem domain and usage of these library functions, human programmers can intuitively decide that the $Lookup$ function, for example, only ever needs to match for one room. Further, they would likely observe that this solution could be made a magnitude faster by using HashMaps instead of a tupled HashSet.

But as programmers deciding on concrete data structures, it is our responsibility to be aware of the restrictions of the concrete types and make the most appropriate decision for the expected data. A compiler should be able to make the same decisions, consulting a longer list of alternatives for further unique cases. If we inform the compiler that HashMaps must represent many-to-one relations in addition to the refined relational types mentioned above (where $location$ and $attends$ are one-to-one and many-to-one respectively), we effectively offload the burden of implementation choice towards a rigid, correct system.

Compare the above solution with the optimal HashMap implementation, which allocates no extra space for intermediate sets and runs in $O(n)$:
\begin{algorithmic}
    \Function{HashMapSolution}{$location: HashMap[Workshop, Room], attends: HashMap[Visitor, Workshop]$}
        \State c := 0
        \For{$v, w \in attends$}
            \If{$location[w] = room$}
                \State $c := c + 1$
            \EndIf
        \EndFor
        \Return $c$
    \EndFunction
\end{algorithmic}

\paragraph{Digression: relational databases} Upon reviewing this example, we see strong parallels in language and functionality to relational database engines (particularly SQL). And as a data processing-focused problem, this example could of course be modelled by a relational database. For a very large event with many rooms, visitors, and workshops, perhaps a relational database is the only way to manage the data. However, using formal notation, the programmer is not burdened with the choice of that representation; they focus on correctness of the abstract data and abstract operations. The specification language intentionally hides implementation-specific operations like joins and counts, instead opting for a purer mathematical model following traditional set theory semantics.

The execution engine should have the freedom of choosing a relational database-like storage-and-query mechanism, but it should also have the freedom to choose between aggressively memory-efficient or fully-cached implementations too. In this particular example, a relational database may not even be the best choice of representation. Imagine we could encode the unchanging carrier sets (ex. a building has a limited number of $Room$s, the set of $Workshop$s is naturally limited by the number of $Room$s) as a compact bitset. If we were to perform an intersection between these encoded sets, the execution could be many times faster than the overhead from a relational database.

By providing the compute engine with more opportunities to follow unique optimization paths, we may be able to tune the computation method to the problem at hand.

% TODO mention collection skeletons area of research - abstract types, with optional traits to provide more domain-related information if known. But the code still works either way

\subsubsection{Programmers should focus on correctness, compiler on efficiency}
One of the best traits of specification languages towards correct program development is their emphasis on proving the intent of the code. Even in industry, organizing software requirements remains a bottleneck for shipping features and products. Although software developers will no doubt need to consider algorithmic efficiency when writing programs, they should be able to abstract away primitive operations to focus on writing readable, correct, and clear code. Then, the compiler should have the freedom to swap out expressive code with a semantically equivalent, more efficient counterpart. In this case, rigorous semantics are required to demonstrate that the source program exhibits the same key behaviours as optimized programs.

High level program libraries and frameworks are already trending towards semantically rich, abstract interfaces covering a low level hand-optimized implementation. For example, the numpy library has become essential syntactic sugar for scientific computing within python. Python further serves as the standard for machine learning tooling, with libraries like PyTorch, TensorFlow, and Triton abstracting complex GPU code behind a vendor-agnostic interface.

By simplifying the language and formalizing semantics, we can grant the compiler more power over optimizations than would otherwise be possible in conventional languages.

\paragraph{Motivating Example: Operation Optimization in Python}
Let's say we have 3 sets $S, T, V$ that can only be represented by the concrete HashSet type. Further, take our candidate query $(S \cup T) \cap V$. Python provides a suite of abstract operators for sets, so this expression may be represented as \texttt{(S | T) \& V}. But its imperative semantics hide the actual computation process:
\begin{minted}{python}
def compute(S, T, V):
    # Step 1: S | T
    s_union_t = set()
    for x in S:
        s_union_t.add(x)
    for x in T:
        s_union_t.add(x)
    # len(s_union_t) <= len(S) + len(T)

    # Step 2: (S | T) & V
    ret = set()
    for x in s_union_t:
        ret.add(x)
    for x in ret:
        if x not in V:
            ret.remove(x)
    # len(ret) <= len(V)
    return ret
\end{minted}
Because the expected behaviour of Python operators is to create a new set while leaving the input arguments untouched, the interpreter must create a new set for the union of $S$ and $T$, then another set for the intersection with $V$. Further, Python is forced to evaluate from left-to-right, obeying the priority of parenthesis. So, in the likely case where $S \cup T$ are larger than $V$, we must waste computation and memory adding and removing values from the return set (as in Step 2). The semantics of conventional imperative languages limit the compiler's ability to shuffle expression evaluation, not least because of side effects.

If we freed the semantics of these operators so that known set theorems can replace expressions, we may choose to swap $(S \cup T) \cap V$ with a more complex but computationally faster expression $(S \cap V) \cup ((T \setminus S) \cap V)$. Further, if we prevent the construction of intermediary sets and instead compute the final result directly, we can create an implementation where the resulting set never undoes prior operations, and never consumes more memory than necessary:
\begin{minted}{python}
def compute(S, T, V):
    # Step 1: S & V
    ret = set()
    for x in S:
        if x in V:
            ret.add(x)
    # len(ret) <= min(len(S), len(V))
    # Step 2
    for x in T:
        if x not in S and x in V:
            ret.add(x)
    # len(ret) <= min(len(S) + len(T), len(V))
    return ret
\end{minted}
The constructed return set strengthens the upper bound cardinality by distributing the intersection operation (which shrinks set sizes) across the union operation (which grows set sizes). As a bonus for focusing on memory efficiency, the runtime is faster too. This second version only iterates over S and T once, while the first version iterates over S, T, and $S \cup T$ twice.

High level specifications can make the intention of computation clearer and simpler, allowing simplifications to happen early in the compilation process. But we must be careful: as demonstrated with this example, the optimized answer is not necessarily the shortest formula. We must consider the concrete type representations, the atomic actions of those concrete types, and the strengths/weaknesses of those atomic actions. A new language defining only abstract semantics allows for extensions of optimizations without any behavioural differences the user would care about; the only behaviours we promise are those defined by the semantics.

% TODO lean into the semantics, proofs, event b literature
% TODO an example of where safety is critical, and language implementations distract from the correctness (or even worse, make something incorrect)

\subsubsection{Partial computation is powerful}
% TODO incrementalization
\paragraph{Motivating Example: Maximum Tail Sum}
%% TODO, benchmarks


% What information is needed for reviewers to understand the design?
% What information is helpful towards contributing to this project?

% Motivation? % Give examples, motivation, why we are taking this route, include steps for code generation, give intuition

\subsection{Syntax}
% Review the syntactic design of simile - scanner/parser arch

% Design decisions:
% - Architecture
% - Generators
% - Iter/Fold

\subsection{Semantics}
% Review the semantic design of simile - how does semantic analysis work? How is it structured? alternatives considered?

% Design decisions:
% - Architecture
% - hilberts choice
% - short circuiting booleans

\subsection{Optimization}
% SimRw - TODO then add SimRW to the optimization specification below

% Design decisions:
% - Architecture
% - SimRW
% - ADT refinement
% - TRS Optimizer + multiple frontends/backends
% - Target? MLIR? LLVM?

A basic strategy to optimize set and relational expressions is:
\begin{enumerate}
  \item Normalize the expression as a set comprehensions
  \item Simplify and reorganize conjuncts of the set comprehension body
\end{enumerate}

The TRS for this language primarily involves lowering collection data type expressions into pointwise boolean quantifications. Breaking down each operation into set builder notation enables a few key actions:
\begin{itemize}
  \item Quantifications over sets ($\bSet{x}{G}{P}$) are naturally separated into generators ($G$) and (non-generating) predicates ($P$). For sets, at least one membership operator per top-level conjunction in $G$ will serve as a concrete element generator in generated code. Then, top level disjunctions will select one membership operation to act as a generator, relegating all others to the predicate level. For example, if the rewrite system observes an intersection of the form $\bSetT{x}{x \in S \land x \in T}$, the set construction operation must iterate over at least one of $S$ and $T$. Then, the other will act as a condition to check every iteration (becoming $\bSet{x}{x \in S}{x \in T}$).
  \item By definition of generators in quantification notation, operations in $G$ must be statements of the form $x \in S$, where $x$ is used in the ``element'' portion of the set construction. Statements like $x \notin T$ or checking a property $p(x)$ must act like conditions since they do not produce any iterable elements.
  \item Any boolean expression for conditions may be rewritten as a combination of $\lnot, \lor$, and $\land$ expressions. Therefore, by converting all set notation down into boolean notation and then generating code based on set constructor booleans, we can accommodate any form of predicate function.
\end{itemize}

% \subsection{Beyond the Language Core}
% Design decisions:
% -

% \paragraph{Granular Strategy (Sets)}
% % TODO take from report.tex, give supporting equations too (all the rules must be stable first)
% \begin{description}
%   \item[Set Comprehension Construction] Break down all qualifying sets into comprehension forms, collapsing and simplifying where needed.
%   \item[DNF Predicates] Revise comprehension predicates to top-level disjunctive normal form. Each or-clause should have at least one feasible generator. Each clause should record a list of candidate generators
%   \item[Predicate Simplification] Remove superfluous dummy variables, group or-clauses that use the exact same generator (ex. $\bSetT{x}{x \in S \land x \neq 0 \lor x \in S \land x = 0} \rightarrow \bSetT{x}{x \in S \land (x \neq 0 \lor x = 0)}$). Clauses should be group-able based on DNF, and generators should be selected and recorded.
%   \item[Set Code Generation] Converts quantifiers into for-loops and if-statements.
% \end{description}

% \section{Supported Operations}
% \begin{table}[H]
%     \centering
%     \caption{Summary table: a few operators on sets and relations.}
%     \begin{tabular}{|c|c||c|c|}
%     \hhline{|--||--|}
%     \multicolumn{2}{|c||}{Sets} & \multicolumn{2}{|c|}{Relations} \\
%     \hhline{:==::==:}
%     Syntax & Label/Description & Syntax & Label/Description\\
%     \hhline{|--||--|}
%     $set(T)$ & Unordered, unique collection             & $S \pfun T$ & Partial function \\
%     $S \leftrightarrow T$ & Relation, $set(S\times T)$  & $S \tinj T$& Total injection\\
%     $\emptyset$ & Empty set                             & $a \mapsto b$ & Pair (relational element) \\
%     $\{a,b,...\}$ & Set enumeration                     & $dom(S)$ & Domain\\
%     $\bSet{x}{x \in S}{P}$ & Set comprehension          & $ran(S)$ & Range\\
%     $S \cup T$ & Union                                  & $R[S]$ & Relational image\\
%     $S \cap T$ & Intersection                           & $R \ovl Q$ & Relational overriding\\
%     $S \setminus T$ & Difference                        & $R \circ Q$ & Relational composition\\
%     $S \times T$ & Cartesian Product                    & $S \triangleleft R$ & Domain restriction\\
%     $S \subseteq T$ & Subset                            & $R^{-1}$ & Relational inverse\\
%     \hhline{|--||--|}
%     % $f(S)$ & Function application\\ % is this like function mapping over a set? do we need to include this?
%     \end{tabular}
%     \label{tab:ADTOps}
% \end{table}
% \begin{table}[H]
%     \centering
%     \caption{Collection of operators on set data types.}
%     \begin{tabular}{|c|c|}
%         \hline
%         Name & Definition \\ %& Type &
%         \hline
%         Empty Set & Creates a set with no elements.\\ %& $\emptyset: set[]$ &
%         Set Enumeration & Literal collection of elements to create a set.\\ %& $\{x, y, ...\}: set[T]$ &
%         Set Membership & The term $x \in S$ is True if $x$ can be found somewhere in $S$. \\ %& $\in: T \times set[T] \rightarrow bool$ &
%         \hline
%         Union & $S \cup T = \bSetT{x}{x \in S \lor x \in T}$ \\ %& $\cup: set[T] \times set[T] \rightarrow set[T]$ &
%         Intersection & $S \cap T = \bSetT{x}{x \in S \land x \in T}$ \\ %& $\cap: set[T] \times set[T] \rightarrow set[T]$ &
%         Difference & $S \setminus T = \bSet{x}{x \in S}{x \notin T}$ \\ %& $\setminus: set[T] \times set[T] \rightarrow set[T]$ &
%         Cartesian Product & $S \times T = \bSetT{x \mapsto y}{x \in S \land y \in T}$ \\ %& $\times: set[T] \times set[V] \rightarrow relation[T,V]$ &
%         \hline
%         Powerset & $\mathbb{P}(S) = \bSetT{s}{s \subseteq S }$ \\ %& $\mathbb{P}: set[T] \rightarrow set[set[T]]$&
%         Magnitude & $\#S = \sum_{x \in S} 1$\\ %& $\#:set[T] \rightarrow int$ &
%         Subset & $S \subseteq T \equiv \forall x \in S: s \in T$ \\ %& $\subseteq: set[T] \times set[T] \rightarrow bool$ &
%         Strict Subset & $S \subset T \equiv S \subseteq T \land S \neq T$ \\ %& $\subset: set[T] \times set[T] \rightarrow bool$ &
%         Superset & $S \supseteq T \equiv \forall x \in T: s \in S$ \\ %& $\supseteq: set[T] \times set[T] \rightarrow bool$ &
%         Strict Superset & $S \supset T \equiv S \supseteq T \land S \neq T$ \\ %& $\supset: set[T] \times set[T] \rightarrow bool$ &
%         \hline
%         Set Mapping & $f * S = \bSetT{f(x)}{x \in S}$\\ %& $*: (T \rightarrow T') \times set[T] \rightarrow set[T']$ &
%         Set Filter & $p \triangleleft S = \bSet{x}{x \in S}{p(x)}$\\ %& $\triangleleft: (T \rightarrow bool) \times set[T] \rightarrow set[T]$ &
%         % Reduction & $f / S = $ %& $/:(T\times T \rightarrow T) \times set[T] \rightarrow T$ &
%         Set Quantification (Folding) & $\oplus x \cdot x \in S \mid P$\\
%         Cardinality & $card(S) = \sum 1 \cdot x \in S$\\
%         \hline
%     \end{tabular}
%     \label{tab:setOps}
% \end{table}
% \begin{table}[H]
%     \centering
%     \caption{Collection of operators on bag/multiset data types.}
%     \begin{tabular}{|c|c|}
%         \hline
%         Name & Definition \\ %& Type &
%         \hline
%         Empty Set & Creates a set with no elements.\\ %& $\emptyset: set[]$ &
%         Bag Enumeration & Literal collection of elements to create a set \\ & (for now, stored as a tuple of elements and number of occurrences).\\ %& $\{x, y, ...\}: set[T]$ &
%         Bag Membership & The term $x \in S$ is True if $S$ contains one or more occurrences of $x$. \\ %& $\in: T \times set[T] \rightarrow bool$ &
%         \hline
%         Union & $S \cup T = \bag{(x, a+b)}{(x,a) \in S \land (x,b) \in T}{a, b \geq 0}$ \\ %& $\cup: set[T] \times set[T] \rightarrow set[T]$ &
%         Intersection & $S \cap T = \bag{(x, min(a,b))}{(x,a) \in S \land (x,b) \in T}{a, b \geq 0}$ \\ %& $\cap: set[T] \times set[T] \rightarrow set[T]$ &
%         Difference & $S - T = \bag{(x, a-b)}{(x,a) \in S \land (x,b) \in T}{a, b \geq 0 \land a-b > 0}$ \\ %& $\setminus: set[T] \times set[T] \rightarrow set[T]$ &
%         % Cartesian Product & $S \times T = \bSetT{x \mapsto y}{x \in S \land y \in T}$ \\ %& $\times: set[T] \times set[V] \rightarrow relation[T,V]$ &
%         % \hline
%         % Powerset & $\mathbb{P}(S) = \bSetT{s}{s \subseteq S }$ \\ %& $\mathbb{P}: set[T] \rightarrow set[set[T]]$&
%         % Magnitude & $\#S = \sum_{x \in S} 1$\\ %& $\#:set[T] \rightarrow int$ &
%         % Subset & $S \subseteq T \equiv \forall x \in S: s \in T$ \\ %& $\subseteq: set[T] \times set[T] \rightarrow bool$ &
%         % Strict Subset & $S \subset T \equiv S \subseteq T \land S \neq T$ \\ %& $\subset: set[T] \times set[T] \rightarrow bool$ &
%         % Superset & $S \supseteq T \equiv \forall x \in T: s \in S$ \\ %& $\supseteq: set[T] \times set[T] \rightarrow bool$ &
%         % Strict Superset & $S \supset T \equiv S \supseteq T \land S \neq T$ \\ %& $\supset: set[T] \times set[T] \rightarrow bool$ &
%         % \hline
%         Bag Mapping & $f * S = \bagT{(f(x), r)}{(x, r) \in S}$\\ %& $*: (T \rightarrow T') \times set[T] \rightarrow set[T']$ &
%         Bag Filter & $p \triangleleft S = \bag{(x, r)}{(x, r) \in S}{p(x)}$\\ %& $\triangleleft: (T \rightarrow bool)
%         Size & $size(S) = \sum r \cdot (x, r) \in S$\\
%         Zero Occurrences & $(x,0) \in S \implies x \notin S$\\
%         \hline
%     \end{tabular}
%     \label{tab:setOps}
% \end{table}
% \begin{table}[H]
%     \centering
%     \caption{Collection of operators on sequence data types.}
%     \begin{tabular}{|c|c|}
%         \hline
%         Name & Definition \\ %& Type &
%         \hline
%         Empty List & Creates a list with no elements.\\ %& $[]: list[]$ &
%         List Enumeration & Literal collection of elements to create a list.\\ %& $[x, y, ...]: list[T]$ &
%         Construction & Alternative form of List Enumeration.\\ %& $(:):T \times list[T] \rightarrow list[T]$ &
%         List Membership & The term $x \texttt{ in } S$ is True if $x$ can be found somewhere in $S$. \\ %& $\texttt{in}: T \times list[T] \rightarrow bool$ &
%         \hline
%         Append & $[s_1, s_2, ..., s_n] + t = [s_1, s_2, ..., s_n, t]$ \\ %& $+: list[T] \times T \rightarrow list[T]$ &
%         Concatenate & $[s_1, ..., s_n] \concat [t_1, ..., t_n] = [s_1, ..., s_n, t_1, ... t_n]$ \\ %& $\concat: list[T] \times list[T] \rightarrow list[T]$ &
%         Length & $\#S = \sum 1 \cdot x \texttt{ in } S$ \\ %& $\#: list[T] \rightarrow int$ &
%         \hline
%         List Mapping & $f * S = \listT{f(x)}{x \texttt{ in } S}$\\ %& $*: (T \rightarrow T') \times list[T] \rightarrow list[T']$ &
%         List Filter & $p \triangleleft S = \list{f(x)}{x \texttt{ in } S}{p(x)}$\\ %& $\triangleleft: (T \rightarrow bool) \times list[T] \rightarrow list[T]$ &
%         Associative Reduction & $\oplus / [s_1, s_2, ..., s_n] = s_1 \oplus s_2 \oplus ... \oplus s_n$\\ %& $/:(T\times T \rightarrow T) \times list[T] \rightarrow T$ &
%         Right Fold & $\texttt{foldr}(f, e, [s_1, s_2, ..., s_n]) = f(s_1 ,f(s_2 , f(..., f(s_n, e))))$\\ %& $\texttt{foldr}:(T\times V \rightarrow V) \times V \times list[T] \rightarrow V$ &
%         Left Fold & $\texttt{foldl}(f, e, [s_1, s_2, ..., s_n]) = f(f(f(f(e, s_1), s_2), ...), s_n)$\\ %& $\texttt{foldl}:(T\times T \rightarrow T) \times list[T] \rightarrow T$ &
%         \hline
%     \end{tabular}
%     \label{tab:seqOps}
% \end{table}
% \begin{table}[H]
%     \centering
%     \caption{Collection of operators on relation data types.}
%     \begin{tabular}{|c|c|}
%         \hline
%         Name & Definition \\ %& Type &
%         \hline
%         Empty Relation & Creates a relation with no elements.\\ %& $\{\}:relation[]$ &
%         Relation Enumeration & Literal collection of elements to create a relation.\\ %& $\{x \mapsto y, a \mapsto b,...\}: relation[T, V]$ &
%         Identity & $id(S)= \bSetT{x \mapsto x}{x \in S}$\\ %& $id: set[T] \rightarrow relation[T,T]$ &
%         Domain & $dom(R) = \bSetT{x}{x \mapsto y \in R}$\\ %& $dom: relation[T,V] \rightarrow set[T]$ &
%         Range & $ran(R) = \bSetT{y}{x \mapsto y \in R}$\\ %& $ran: relation[T,V] \rightarrow set[V]$ &
%         \hline
%         Relational Image & $R[S] = \bSet{y}{x \mapsto y \in R}{x \in S}$ \\ %& $([]): relation[T,V] \times set[T] \rightarrow set[V]$ &
%         Overriding & $R \ovl Q = Q \cup (dom(Q) \domsub R)$\\ %\Set{x \mapsto y}{x\mapsto y \in Q \lor (x \mapsto y \in R \land x \notin dom(Q))}$ \\ %& $\ovl: relation[T,V] \times relation[T,V] \rightarrow relation[T,V]$ &
%         (Forward) Composition & $Q \circ R = \bSetT{x \mapsto z}{x \mapsto y \in R \land y \mapsto z \in Q}$\\ %& $\circ: relation[V,W] \times relation[T,V] \rightarrow relation[T,W]$ &
%         Inverse & $R^{-1} = \bSetT{y \mapsto x}{x \mapsto y \in R}$ \\ %& $(^{-1}): relation[T,V] \rightarrow relation[V,T]$ &
%         \hline
%         Domain Restriction & $S \triangleleft R = \bSet{x \mapsto y}{x \mapsto y \in R}{x \in S}$\\ %& $\triangleleft: set[T] \times relation[T,V] \rightarrow relation[T,V]$ &
%         Domain Subtraction & $S \domsub R = \bSet{x \mapsto y}{x \mapsto y \in R}{x \notin S}$\\ %& $\domsub: set[T] \times relation[T,V] \rightarrow relation[T,V]$ &
%         Range Restriction & $R \triangleright S = \bSet{x \mapsto y}{x \mapsto y \in R}{y \in S}$\\ %& $\triangleright: set[V] \times relation[T,V] \rightarrow relation[T,V]$ &
%         Range Subtraction & $R \ransub S = \bSet{x \mapsto y}{x \mapsto y \in R}{y \notin S}$\\ %& $\ransub: set[V] \times relation[T,V] \rightarrow relation[T,V]$ &
%         % \hline % Do we need to include tests for total, surjective, injective, etc.?
%         % Eventually add closures
%         \hline
%     \end{tabular}
%     \label{tab:relOps}
% \end{table}

\section{Specification}
This section contains the complete specification and low-level design of Simile, inspired by Event-B and Python.

\subsection{Lexicon}
Since Simile makes extensive use of Event-B-like set notation and otherwise non-ASCII characters, this section contains a list of such characters, translations to ASCII symbols where suitable, and other useful lexing information.

\subsubsection{Whitespace}
Similar to Python, indentation and newline tokens are significant in Simile. In particular, nested code blocks, as in if-statements, for-loops, and procedure definitions, require the use of indentations to signify the beginning and end. Indentation/newlines may also be used within collection-type comprehensions or enumerations, similar to the style seen in JSON. Comments are single-lined and start with a pound symbol (\#). Indentation is ignored within bracketed expressions\footnote{Except for within the `iter' block. This is a special case handled post-lexing: ignore indentation when parenthesis/bracket/brace level $>1$ and we are not within the block portion of `iter'... VBAR NEWLINE INDENT `block' DEDENT.}.

\subsubsection{Identifiers}
Identifiers in Simile follow the standard regular expression, using only ASCII letters, underscores, and numbers: $[a-zA-Z\_][a-zA-Z\_0-9]$.

\subsubsection{Primitives}
Integers and floating point numbers follow a similar style to modern programming languages. The dash character ($-$) may be used in either subtraction or as a unary negation. Superscript numbers may be used for constant exponentiation. Booleans (True/False) follow Python style, with capitalized first letters. Strings are delimited with double-quotes (") and may be multilined.

% TODO: escaped characters in strings?

\subsubsection{Operators}
ASCII tokens will be provided as an alternative to unicode format where suitable. The following are symbols used for expressions:
\begin{center}
\begin{longtable}{c c l}
    \textbf{Token} & \textbf{ASCII Token} & \textbf{Name} \\
    \endhead
    \lexentry{\cdot}{.:}{Center Dot (Quantifier Separator) \footnote{Anywhere $\cdot$ is valid, $.$ is valid, except the parser may attempt to treat $.$ as a record field access instead of a quantifier separator (like $a.b$), so we prefer the use of $\cdot$ instead.}} % TODO should we not allow this, and instead just stick with .:?
    \lexentry{.}{}{Dot}
    \lexentry{,}{iter\_and}{Comma}
    \lexentry{`}{iter\_or}{Comma}
    \lexentry{:}{}{Colon}
    \lexentry{;}{}{Semicolon}
    \lexentry{|}{}{Vertical Bar}
    \lexentry{\lambda}{lambda}{Lambda}
    \lexentry{(}{}{Left Parenthesis}
    \lexentry{)}{}{Right Parenthesis}
    \lexentry{\langle}{<{}<}{Left Bracket}
    \lexentry{}{[}{Left Bracket (Alternative)}
    \lexentry{\rangle}{>{}>}{Right Bracket}
    \lexentry{}{]}{Right Bracket (Alternative)}
    \lexentry{\{}{}{Left Brace}
    \lexentry{\}}{}{Right Brace}
    \lexentry{\llbracket}{[[}{Left Double Bracket}
    \lexentry{\rrbracket}{]]}{Right Double Bracket}
    \lexentry{=}{}{Equals}
    \lexentry{\neq}{!=}{Not Equals}
    \lexentry{\lnot}{not}{Not}
    \lexentry{}{!}{Not (Alternative)}
    \lexentry{\land}{and}{And}
    \lexentry{\lor}{or}{Or}
    \lexentry{\Rightarrow}{=>}{Implies}
    % \lexentry{\Leftarrow}{<==}{Reverse Implies}
    \lexentry{\equiv}{==}{Equivalent}
    % \lexentry{\equiv}{}{Equivalent (Alternative)}
    \lexentry{\not\equiv}{!==}{Not Equivalent}
    \lexentry{\forall}{forall}{For All}
    \lexentry{\exists}{exists}{There Exists}
    \lexentry{}{+}{Plus}
    \lexentry{}{-}{Minus}
    \lexentry{}{*}{Multiplication}
    \lexentry{}{/}{Division}
    \lexentry{}{div}{Integer Division}
    \lexentry{}{mod}{Modulo}
    \lexentry{}{\^}{Exponent}
    \lexentry{<}{}{Less Than}
    \lexentry{\leq}{<=}{Less Than or Equal}
    \lexentry{>}{}{Greater Than}
    \lexentry{\geq}{>=}{Greater Than or Equal}
    \lexentry{\in}{in}{In (Membership)}
    \lexentry{\not\in}{not in}{Not In (Membership)}
    \lexentry{\cup}{\textbackslash/}{Union}
    \lexentry{\cap}{/\textbackslash}{Intersection}
    \lexentry{\setminus}{\textbackslash}{Backslash (Set Minus)}
    \lexentry{\times}{><}{Cartesian Product}
    % \lexentry{\times}{><}{Cartesian Product}
    \lexentry{\subset}{<{}<:}{Subset}
    \lexentry{\subseteq}{<:}{Subset or Equal}
    \lexentry{\supset}{:>{}>}{Superset}
    \lexentry{\supseteq}{:>}{Superset or Equal}
    \lexentry{\not\subset}{!<{}<:}{Not Subset}
    \lexentry{\not\subseteq}{!<:}{Not Subset or Equal}
    \lexentry{\not\supset}{!:>{}>}{Not Superset}
    \lexentry{\not\supseteq}{!:>}{Not Superset or Equal}
    \lexentry{\bigcup}{union}{General Union}
    \lexentry{\bigcap}{intersection}{General Intersection}
    \lexentry{\mapsto}{|->}{Maplet}
    \lexentry{\oplus}{<+>}{Relation Overriding} % Old: \ovl \footnote{This unicode symbol requires a specific font available from Rodin: Brave Sans Mono, U+E103}}
    \lexentry{\circ}{<>}{Composition}
    \lexentry{\concat}{++}{Concatenation}
    \lexentry{^{-1}}{\~}{Inverse}
    \lexentry{\domres}{<|}{Domain Restriction}
    \lexentry{\domsub}{<{}<|}{Domain Subtraction}
    \lexentry{\ranres}{|>}{Range Restriction}
    \lexentry{\ransub}{|>{}>}{Range Subtraction}
    \lexentry{\rel}{<->}{Relation}
    \lexentry{\trel}{<{}<->}{Total Relation \footnote{Brave Sans Mono unicode: U+E100}}
    \lexentry{\srel}{<->{}>}{Surjective Relation \footnote{Brave Sans Mono unicode: U+E101}}
    \lexentry{\strel}{<{}<->{}>}{Total Surjective Relation \footnote{Brave Sans Mono unicode: U+E102}}
    \lexentry{\pfun}{+->}{Partial Function}
    \lexentry{\tfun}{-{}->}{Total Function}
    \lexentry{\pinj}{>+>}{Partial Injection}
    \lexentry{\tinj}{>->}{Total Injection}
    \lexentry{\psur}{+->{}>}{Partial Surjection}
    \lexentry{\tsur}{-{}->{}>}{Total Surjection}
    \lexentry{\tbij}{>->{}>}{Bijection}
    \lexentry{\upto}{}{Upto}
    \lexentry{\pow}{powerset}{Powerset}
    \lexentry{\pown}{powerset1}{Non-Empty Powerset}
    \lexentry{\sum}{sum}{Summation}
    \lexentry{\prod}{product}{Product}
\end{longtable}
\end{center}

The following are symbols used for programming constructs:
\begin{center}
\begin{longtable}{c c c c}
    \textbf{(Unicode) Token} & \textbf{(ASCII) Token} & \textbf{Name} \\
    \endhead
    \lexentry{\coloneqq}{:=}{Assignment}
    \lexentry{:\in}{::}{Choice (Non-deterministic Membership) Assignment}
    \lexentry{\rightarrow}{->}{Right Arrow}
\end{longtable}
\end{center}

\subsubsection{Reserved Keywords}
The following list of symbols and identifiers are reserved for programming constructs:

\begin{center}
\begin{tabular}{c c c c c c}
true & false & if & else & for\\
while & record & enum & procedure & is\\
is not & return & break & continue & from\\
import & skip & trait & type & iter \\
iter\_or & iter\_and
\end{tabular}
\end{center}

And for built-in types:
\begin{center}
\begin{tabular}{c c c c c c}
int &
float &
string &
bool &
$\mathbb{Z}$ &
$\mathbb{N}$ \\
$\mathbb{N}_1$ &
set &
sequence &
bag &
relation &
generic \\
tuple &
\end{tabular}
\end{center}

And built-in procedures:

\begin{center}
\begin{tabular}{c c c c c c}
min &
map\_min &
max &
map\_max &
choice &
dom \\
ran &
card &
size &
sum &
cast &
cast\_with \\
print &
map &
heads &
tails &
fold
\end{tabular}
\end{center}

\subsection{Syntax}
We follow the standard EBNF notation for describing the language structure.

\begin{grammar}
<Start> ::= <Statements>

<Statements> ::= \{ <SimpleStatements> | <CompoundStmt>\}

% Small hack needed with trait statements hijacking assignment or expr since lots of productions would share the `first' set
<SimpleStatements> ::= <AssignmentOrExpr> (`;' <SimpleStmt> \{`;' <SimpleStmt>\} NEWLINE | NEWLINE [<TraitStmt>])
    \alt (<ControlFlowStmt> | <ImportStmt>) \{`;' <SimpleStmt>\} NEWLINE

<SimpleStmt> ::= <AssignmentOrExpr>
    \alt <ControlFlowStmt>
    \alt <ImportStmt>

<AssignmentOrExpr> ::= <Expr> [`:' <TypeExpr>] [(\lit{$\coloneqq$} | \lit{$:\in$}) <Expr>]

<TraitStmt> ::= INDENT `trait' <Expr> NEWLINE \{`trait' <Expr> NEWLINE\} DEDENT

<Expr> ::= <Quantification> | <Predicate>

<Quantification> ::= \lit{$\lambda$} [<FlatTupleIdentifier>] \lit{$\cdot$} <Predicate> `|' <Expr>
    \alt (\lit{$\sum$} | \lit{$\prod$} | \lit{$\bigcup$} | \lit{$\bigcap$} | \lit{$\exists$} | \lit{$\forall$} | `max' | `min') <QuantificationBody>
    % fold requires a generator, an accumulator assignment with an initial value, and an update expr (no need to explicitly assign to the accumulator var)
    \alt `fold' <Assignment> `:' <QuantificationBody>
    % iter requires a generator, initialization values, maybe a return value/identifier/tupleIdentifier, then some statements to modify the init values
    \alt `iter' <IterQuantificationBody>

<QuantificationBody> ::= <BranchQuantificationBody>
    % Generator trees all end in expr eventually (even if that expr is nested within branches)
    \alt <Generator> \{`,' <Generator>\} (`,' <BranchQuantificationBody> | `|' <Expr>)

<BranchQuantificationBody> ::= `(' <QuantificationBody> `)' \{``(' <QuantificationBody> `)'\}

% Predicates are attached to a single generator and only apply to that generator
<Generator> ::= <IdentList> \lit{$\in$} <Expr> [`·' <Expr>]

<IterQuantificationBody> ::= <BranchIterQuantificationBody>
    \alt <GeneratorWithAssignments> \{`,' <GeneratorWithAssignments> \} (`,' <BranchIterQuantificationBody> | `->' <IdentList> `|' <IterBlock>)

<BranchIterQuantificationBody> ::= `(' <IterQuantificationBody> `)' \{``(' <IterQuantificationBody> `)' \}

<GeneratorWithAssignments> ::= <Generator> `:' <Assignment> \{ `;' <Assignment>\}

<Assignment> ::= <Expr> [`:' <TypeExpr>] (\lit{$\coloneqq$} | \lit{$:\in$}) <Expr>

% Iterblock separate from <Block> so that we can safely ignore indentation when we do not see the VBAR NEWLINE INDENT combo
<IterBlock> ::= <SimpleStmt> \{`;' <SimpleStmt>\}
    \alt NEWLINE INDENT <Statements> DEDENT

<IdentList> ::= <IdentListItem> \{`,' <IdentListItem>\}

<IdentListItem> ::= IDENTIFIER
    \alt <IdentListItem> \lit{$\mapsto$} <IdentListItem>
    \alt `(' <IdentList> `)'

<Predicate> ::= <Implication>
    \alt <Predicate> (\lit{$\equiv$} | \lit{$\not\equiv$}) <Implication>

% TODO is this right? rev implies seems awkward
<Implication> ::= <Disjunction>
    \alt <Disjunction> \{\lit{$\Rightarrow$} <Disjunction>\}
    % \alt <Disjunction> \{\lit{$\impliedby$} <Implication>\}

<Disjunction> ::= <Conjunction> \{\lit{$\lor$} <Conjunction>\}

<Conjunction> ::= <Negation> \{\lit{$\land$} <Negation>\}

<Negation> ::= <AtomBool>
    \alt \lit{$\lnot$} <Negation>

<AtomBool> ::= <PairExpr> [(
        \lit{$=$}
      | \lit{$\neq$}
      % | \lit{is}
      % | \lit{is not}
      | `<'
      | `>'
      | \lit{$\leq$}
      | \lit{$\geq$}
      | \lit{$\in$}
      | \lit{$\not\in$}
      | \lit{$\subset$}
      | \lit{$\subseteq$}
      | \lit{$\supset$}
      | \lit{$\supseteq$}
      | \lit{$\not\subset$}
      | \lit{$\not\subseteq$}
      | \lit{$\not\supset$}
      | \lit{$\not\supseteq$}
    ) <PairExpr>]

<PairExpr> ::= <RelSetExpr>
    \alt <PairExpr> \lit{$\mapsto$} <RelSetExpr>

<RelSetExpr> ::= <SetExpr> [(
      \lit{$\rel$}
    | \lit{$\trel$}
    | \lit{$\srel$}
    | \lit{$\strel$}
    | \lit{$\pfun$}
    | \lit{$\tfun$}
    | \lit{$\pinj$}
    | \lit{$\tinj$}
    | \lit{$\psur$}
    | \lit{$\tsur$}
    | \lit{$\tbij$}
) <RelSetExpr>]

<SetExpr> ::= <IntervalExpr>
    \alt <IntervalExpr> \{\lit{$\cup$} <IntervalExpr>\}
    \alt <IntervalExpr> \{\lit{$\times$} <IntervalExpr>\}
    \alt <IntervalExpr> \{\lit{$\oplus$} <IntervalExpr>\}
    \alt <IntervalExpr> \{\lit{$\circ$} <IntervalExpr>\}
    \alt <IntervalExpr> \{\lit{$\cap$} <IntervalExpr>\}
    \alt <IntervalExpr> \{\lit{$\concat$} <IntervalExpr>\}
    \alt <IntervalExpr> (\lit{$\domsub$} | \lit{$\domres$}) <IntervalExpr> \{\lit{$\cap$} <IntervalExpr>\} [<RelSubExpr>]

<RelSubExpr> ::= (\lit{$\setminus$} | \lit{$\ranres$} | \lit{$\ransub$}) <IntervalExpr>

<IntervalExpr> ::= <ArithmeticExpr> [`..' <ArithmeticExpr>]

<ArithmeticExpr> ::= <Term>
    \alt <ArithmeticExpr> (`+' | `-') <Term>

% Need to have \times and * be different symbols - otherwise 1..n \times n is ambiguous - either the set of 1 to n*n or the set of 1..n, n times (I know the semantics for the second one dont make much sense, but it may be best to separate the two symbols because of the differing precedence around .. and +/-
<Term> ::= <Factor>
    \alt <Term> (`*' | `/' | `div' | `mod') <Factor>

<Factor> ::= [(`+' | `-')] <Power>

% Note, because of next production, if you see {}^{-1}, look for relational image operation
<Power> ::= <Primary> [\lit{\^{}} <Factor> | [\lit{${}^-$}](\lit{${}^{INTEGER}$} | \lit{${}^{FLOAT}$})]

<Primary> ::= <Atom> \{<AtomFollow>\}

<AtomFollow> ::= `.' IDENTIFIER
    \alt `(' <Expr> \{`,' <Expr>\} `)'
    % comma within this image operation is to allow for parameterized generics similar to python's typing. Maybe we should extract into a different object?
    \alt `[' <Expr> \{`,' <Expr> \} `]' % Python like list slcing but with commas? or should we allow multiple things to be queried... Primary use of commas is for type exprs

<Atom> ::= INTEGER | FLOAT | STRING | BOOLEAN | IDENTIFIER
    \alt <Set> | <Sequence> | <Bag> | <Tuple>
    %\alt \lit{$\mathbb{P}$} `(' <Expr> `)' | \lit{$\mathbb{P}_1$} `(' <Expr> `)'
    \alt `(' <Expr> `)' % SEE NOTE for Tuple

<Set> ::= `{' <CollectionBody> `}'

<Bag> ::= \lit{$\llbracket$} <CollectionBody> \lit{$\rrbracket$}

<Sequence> ::= \lit{$\langle$} <CollectionBody> \lit{$\rangle$}

% Tuples either are empty, have one element (distinct from a parenthesized expr by a comma), or follow enumerationbody
% TODO check if this attempts to match to parenth expr in atom... may need to move tuple up to that function in the code
<Tuple> ::= `(' `)' | `(' [NEWLINE INDENT] <Expr> `,' [NEWLINE [INDENT]] [<EnumerationBody>] `)'

<CollectionBody> ::= <QuantificationBody>
    \alt <EnumerationBody>

<EnumerationBody> ::= [NEWLINE INDENT] <Expr> \{`,' [NEWLINE [INDENT]] <Expr>\} [NEWLINE DEDENT]

<CompoundStmt> ::= <IfStmt> | <ForStmt> | <WhileStmt>
    \alt <RecordStmt> | <ProcedureStmt>

<IfStmt> ::= `if' <Predicate> `:' <Block> [<ElseStmt>]

<ElseStmt> ::= `else :' <Block>
    \alt `else if' <Predicate> `:' <Block> [<ElseStmt>]

<ForStmt> ::= `for' <IdentList> \lit{$\in$} <Expr> `:' <Block>

<WhileStmt> ::= `while' <Expr> `:' <Block>

<RecordStmt> ::= `record' IDENTIFIER `:' NEWLINE INDENT ( \\
        `skip' | <TypedName> \{(`,' [NEWLINE] | NEWLINE) <TypedName> \} \\
    ) NEWLINE DEDENT

<ProcedureStmt> ::= `procedure' IDENTIFIER `(' <TypedName> \{`,' <TypedName> \} `)' \lit{$\rightarrow$} <Expr> `:' <Block>

<Block> ::= <SimpleStatements>
    \alt NEWLINE INDENT <Statements> DEDENT

<TypedName> ::= IDENTIFIER [`:' <TypeExpr>]

<TypeExpr> ::= IDENTIFIER \{`.' IDENTIFIER\} [`[' <TypeExpr> \{`,' <TypeExpr>\}`]']
    \alt `(' `)'
    \alt `(' <TypeExpr> `,' [<TupleTypeExprBody>] `)'

<TupleTypeExprBody> ::= <TypeExpr> \{`,' <TypeExpr>\}

<ControlFlowStmt> ::= `return' [<Expr>]
    \alt `break'
    \alt `continue'
    \alt `skip'

<ImportStmt> ::= `import' STRING
    \alt `from' STRING `import' <ImportList>

<ImportList> ::= `*'
    \alt <FlatTupleIdentifier>

<FlatTupleIdentifier> ::= IDENTIFIER \{`,' IDENTIFIER \}
    \alt `(' IDENTIFIER \{`,' IDENTIFIER \} `)'

\end{grammar}


\subsection{Static Analysis - Type System}
Let $\Gamma$ be the type environment set consisting of elements either $identifier: Type$ or $Type$. The typing rules have been organized in distinct categories for readability. The notation $Type[x, y, z]$ denotes refinements on $Type$ either through instantiating parametric (generic) types or restricting values by way of \texttt{trait} clauses.

Type refinements may be applicable to types with certain traits (ex., min/max values for ordered types, definedness of expressions with certain values, set sizes). % Is there a way to formalize this? Refinement calculus?

\subsubsection{Base Environment Interactions}
\paragraph{Includes:} Mechanisms for adding/extracting variables from the type environment.

\begin{gather}
\tag{Env $\emptyset$}
\begin{prooftree}
\infer0{\emptyset \vdash \diamond}
\end{prooftree}
\\
\tag{Env $I$}
\begin{prooftree}
\hypo{\Gamma \vdash A}
\hypo{I \not\in dom(\Gamma)}
\infer2{\Gamma, I:A \vdash \diamond}
\end{prooftree}
\\
\tag{Fetch Identifier}
\begin{prooftree}
\hypo{\Gamma_1, I:A, \Gamma_2 \vdash \diamond}
\infer1{\Gamma_1, I:A, \Gamma_2 \vdash I:A}
\end{prooftree}
\end{gather}

\subsubsection{Subtype Interactions}
\paragraph{Includes:} Generic subtyping rules.

\begin{gather}
\tag{Reflexive Subtype}
\begin{prooftree}
\hypo{\Gamma \vdash A}
\infer1{\Gamma \vdash A \leq A}
\end{prooftree}
\\
\tag{Transitive Subtype}
\begin{prooftree}
\hypo{\Gamma \vdash A \leq B}
\hypo{\Gamma \vdash B \leq C}
\infer2{\Gamma \vdash A \leq C}
\end{prooftree}
\\
\tag{Subsumption}
\begin{prooftree}
\hypo{\Gamma \vdash a:A}
\hypo{\Gamma \vdash A \leq B}
\infer2{\Gamma \vdash a: B}
\end{prooftree}
\\
\tag{Top Type}
\begin{prooftree}
\hypo{\Gamma \vdash \diamond}
\infer1{\Gamma \vdash Any}
\end{prooftree}
\\
\tag{Sub Top Type}
\begin{prooftree}
\hypo{\Gamma \vdash A}
\infer1{\Gamma \vdash A \leq Any}
\end{prooftree}
\\
\tag{Sub Function}
\begin{prooftree}
\hypo{\Gamma \vdash A' \leq A}
\hypo{\Gamma \vdash B \leq B'}
\infer2{\Gamma \vdash A \rightarrow B \leq A' \rightarrow B'}
\end{prooftree}
\\
\tag{Sub Set}
% TODO check this - doesnt match up with the definition of set as set: A -> bool, but does work intuitively
% ex. set[nat] <= set[int]
\begin{prooftree}
\hypo{\Gamma \vdash A \leq B}
\infer1{\Gamma \vdash set[A] \leq set[B]}
\end{prooftree}
\\
\tag{Sub Product}
\begin{prooftree}
\hypo{\Gamma \vdash A_1 \leq B_1}
\hypo{\Gamma \vdash A_2 \leq B_2}
\infer2{\Gamma \vdash A_1 \times A_2 \leq B_1 \times B_2 }
\end{prooftree}
\\
\tag{Sub Record}
\begin{prooftree}
\hypo{\Gamma \vdash A_1 \leq B_1 ... A_n \leq B_n}
\hypo{\Gamma \vdash A_{n+1} ... A_{n+m}}
\infer2{\Gamma \vdash  Record[I_1: A_1, ..., I_n: A_n, I_{n+1}: A_{n+1}, ... I_{n+m}: A_{n+m}] \leq  Record[I_1: B_1, ..., I_n: B_n]}
\end{prooftree}
\\
\tag{Type Refinement}
\begin{prooftree}
\hypo{\Gamma \vdash T}
\infer1{\Gamma \vdash T[\text{with } x] \leq T}
\end{prooftree}
\end{gather}

\subsubsection{The Set Type}
\paragraph{Includes:} Universal type; everything is a set. The empty set is a part of every type (aside from primitives).

\begin{gather}
\tag{Powerset}
\begin{prooftree}
\hypo{\Gamma \vdash T}
\infer1{\Gamma \vdash set[T]}
\end{prooftree}
\\
\tag{Emptyset}
\begin{prooftree}
\hypo{\Gamma \vdash \diamond}
\infer1{\Gamma \vdash EmptySet}
\end{prooftree}
\\
\tag{Emptyset Bottom}
\begin{prooftree}
\hypo{\Gamma \vdash EmptySet, set[T]}
\infer1{EmptySet \leq set[T]}
\end{prooftree}
\\
\tag{Set Enumeration}
\begin{prooftree}
\hypo{\Gamma \vdash e_1, ..., e_n: A}
\infer1{\Gamma \vdash \{e_1, ..., e_n\}: set[A]}
\end{prooftree}
\end{gather}

\subsubsection{Primitive Types}
\paragraph{Includes:} Basic types built in to the language, like numbers (int and float), strings, and booleans. Although these could all be represented as sets, it is useful to distinguish these common types from other collections or user-defined types.

\begin{gather}
\tag{Primitives - bool}
\begin{prooftree}
\hypo{\Gamma \vdash \diamond}
\infer1{\Gamma \vdash bool }
\end{prooftree}
\\
\tag{Primitives - int}
\begin{prooftree}
\hypo{\Gamma \vdash \diamond}
\infer1{\Gamma \vdash int }
\end{prooftree}
\\
\tag{Primitives - float}
\begin{prooftree}
\hypo{\Gamma \vdash \diamond}
\infer1{\Gamma \vdash float }
\end{prooftree}
\\
\tag{Primitives - str}
\begin{prooftree}
\hypo{\Gamma \vdash \diamond}
\infer1{\Gamma \vdash str }
\end{prooftree}
\\
\tag{Nat from int}
\begin{prooftree}
\hypo{\Gamma \vdash int}
\infer1{\Gamma \vdash nat \leq int[\text{with } min=0] }
\end{prooftree}
\\
\tag{Int from float}
\begin{prooftree}
\hypo{\Gamma \vdash int, float}
\infer1{int \leq float }
\end{prooftree}
\end{gather}

\subsubsection{Collection Types}
\paragraph{Includes:} Sequences, relations, enums (frozen sets), and bags (multisets); Syntactic sugar for abstract data types defined in terms of sets.

\begin{gather}
\tag{Tuples}
\begin{prooftree}
\hypo{\Gamma \vdash T_1, ..., T_n}
\infer1{\Gamma \vdash tuple[T_1, ..., T_n] \quad tuple[T_1, ..., T_n] = T_1 \times ... \times T_n}
\end{prooftree}
\\
\tag{Empty Tuple}
\begin{prooftree}
\hypo{\Gamma \vdash \diamond}
\infer1{\Gamma \vdash (): EmptyTuple}
\end{prooftree}
\\
\tag{Empty Tuple Bottom}
\begin{prooftree}
\hypo{\Gamma \vdash EmptyTuple, T}
\infer1{\Gamma \vdash EmptyTuple \leq tuple[T]}
\end{prooftree}
\\
\tag{Tuple Enumeration}
\begin{prooftree}
\hypo{\Gamma \vdash e_1:A_1, ..., e_n: A_n}
\infer1{\Gamma \vdash ( e_1, ..., e_n ) : tuple[A_1, ..., A_n]}
\end{prooftree}
\\
\tag{Relation Type}
\begin{prooftree}
\hypo{\Gamma \vdash set[T_1], set[T_2]}
\infer1{\Gamma \vdash T_1 \rel T_2 \quad T_1 \rel T_2 = set[tuple[T_1, T_2]]}
\end{prooftree}
\\
\tag{Bag Type}
\begin{prooftree}
\hypo{\Gamma \vdash T, nat}
\infer1{\Gamma \vdash bag[T] \leq T \pfun nat }
\end{prooftree}
\\
\tag{Bag Enumeration}
\begin{prooftree}
\hypo{\Gamma \vdash e_1, ..., e_n: A}
\infer1{\Gamma \vdash \llbracket e_1, ..., e_n \rrbracket: bag[A]}
\end{prooftree}
\\
\tag{Sequence Type}
\begin{prooftree}
\hypo{\Gamma \vdash T, nat}
\infer1{\Gamma \vdash sequence[T] \leq nat \pfun T }
\end{prooftree}
\\
\tag{Sequence Enumeration}
\begin{prooftree}
\hypo{\Gamma \vdash e_1, ..., e_n: A}
\infer1{\Gamma \vdash \langle e_1, ..., e_n \rangle : sequence[A]}
\end{prooftree}
\\
\tag{Enum from static set}
\begin{prooftree}
\hypo{\Gamma \vdash s_1, ..., s_n : str}
\hypo{A = \{s_1, ..., s_n\} }
\hypo{immutable(A)}
\infer3{\Gamma \vdash enum[A] \leq set[A, \text{with } immutable]}
\end{prooftree}
\end{gather}



\subsubsection{Relation Subtypes}
\paragraph{Includes:} Relation subtypes examine four key properties of relations: totality ($t$), surjectivity ($t_r$), one-to-manyness ($1-$), and many-to-oneness ($1-_r$). Syntactic sugar to produce refined relations with these properties follows the definitions presented in Event-B. See Section 5.1.3 for more details.

\begin{gather}
\tag{Relation Subtype - Total Relation}
\begin{prooftree}
\hypo{\Gamma \vdash set[T_1], set[T_2]}
\infer1{\Gamma \vdash T_1 \trel T_2 \quad T_1 \trel T_2 = (T_1 \rel T_2)[\text{with } t]}
\end{prooftree}
\\
\tag{Relation Subtype - Surjective Relation}
\begin{prooftree}
\hypo{\Gamma \vdash set[T_1], set[T_2]}
\infer1{\Gamma \vdash T_1 \srel T_2 \quad T_1 \srel T_2 = (T_1 \rel T_2)[\text{with } t_r]}
\end{prooftree}
\\
\tag{Relation Subtype - Total Surjective Relation}
\begin{prooftree}
\hypo{\Gamma \vdash set[T_1], set[T_2]}
\infer1{\Gamma \vdash T_1 \strel T_2 \quad T_1 \strel T_2 = (T_1 \rel T_2)[\text{with } t, \text{ with } t_r]}
\end{prooftree}
\\
\tag{Relation Subtype - Partial Function}
\begin{prooftree}
\hypo{\Gamma \vdash set[T_1], set[T_2]}
\infer1{\Gamma \vdash T_1 \pfun T_2 \quad T_1 \pfun T_2 = (T_1 \rel T_2)[\text{with } 1-]}
\end{prooftree}
\\
\tag{Relation Subtype - Total Function}
\begin{prooftree}
\hypo{\Gamma \vdash set[T_1], set[T_2]}
\infer1{\Gamma \vdash T_1 \tfun T_2 \quad T_1 \tfun T_2 = (T_1 \rel T_2)[\text{with } t, \text{ with } 1-]}
\end{prooftree}
\\
\tag{Relation Subtype - Partial Injection}
\begin{prooftree}
\hypo{\Gamma \vdash set[T_1], set[T_2]}
\infer1{\Gamma \vdash T_1 \pinj T_2 \quad T_1 \pinj T_2 = (T_1 \rel T_2)[\text{with } 1-, \text{ with } 1-_r]}
\end{prooftree}
\\
\tag{Relation Subtype - Total Injection}
\begin{prooftree}
\hypo{\Gamma \vdash set[T_1], set[T_2]}
\infer1{\Gamma \vdash T_1 \tinj T_2 \quad T_1 \tinj T_2 = (T_1 \rel T_2)[\text{with } t, \text{ with } 1-, \text{ with } 1-_r]}
\end{prooftree}
\\
\tag{Relation Subtype - Partial Surjection}
\begin{prooftree}
\hypo{\Gamma \vdash set[T_1], set[T_2]}
\infer1{\Gamma \vdash T_1 \psur T_2 \quad T_1 \psur T_2 = (T_1 \rel T_2)[\text{with } t_r, \text{ with } 1-]}
\end{prooftree}
\\
\tag{Relation Subtype - Total Surjection}
\begin{prooftree}
\hypo{\Gamma \vdash set[T_1], set[T_2]}
\infer1{\Gamma \vdash T_1 \tsur T_2 \quad T_1 \tsur T_2 = (T_1 \rel T_2)[\text{with } t, \text{ with } t_r, \text{ with } 1-]}
\end{prooftree}
\\
\tag{Relation Subtype - Bijection}
\begin{prooftree}
\hypo{\Gamma \vdash set[T_1], set[T_2]}
\infer1{\Gamma \vdash T_1 \tbij T_2 \quad T_1 \tbij T_2 = (T_1 \rel T_2)[\text{with } t, \text{ with } t_r, \text{ with } 1-, \text{ with } 1-_r]}
\end{prooftree}
\end{gather}

\subsubsection{Simple Commands}
\paragraph{Includes:} Assignment, type assignment, type refinements, simple programming language commands/statements.

\begin{gather}
\tag{Variable Assignment}
\begin{prooftree}
\hypo{\Gamma \vdash E:A}
\hypo{I \not\in dom(\Gamma) \lor \Gamma \vdash I:A}
\infer2{\Gamma \vdash I:A := E }
\end{prooftree}
\\
\tag{Type Alias Assignment}
\begin{prooftree}
\hypo{\Gamma \vdash T}
\hypo{I \not\in dom(\Gamma)}
\infer2{\Gamma \vdash I := T }
\end{prooftree}
\\
\tag{Type Alias}
\begin{prooftree}
\hypo{\Gamma \vdash I := T}
\infer1{\Gamma, I \vdash \diamond \quad I = T}
\end{prooftree}
\\
\tag{Refined Variable Assignment}
\begin{prooftree}
\hypo{\Gamma \vdash E: A, E_1: A_1, ..., E_n:A_n}
\hypo{(I \not\in dom(\Gamma) \lor \Gamma \vdash I:A)}
\infer2{\Gamma \vdash I: A[\text{with } E_1: A_1, ..., \text{ with } E_n: A_n] := E}
\end{prooftree}
\\
\tag{Refined Type Alias}
\begin{prooftree}
\hypo{\Gamma \vdash T, E_1: A_1, ..., E_n: A_n}
\hypo{I \not\in dom(\Gamma)}
\infer2{\Gamma \vdash I := T[\text{with } E_1: A_1, ...,\text{ with } E_n: A_n]}
\end{prooftree}
\\
\tag{Command - break}
\begin{prooftree}
\hypo{\Gamma \vdash \diamond}
\infer1{\Gamma \vdash \text{break}: C}
\end{prooftree}
\\
\tag{Command - continue}
\begin{prooftree}
\hypo{\Gamma \vdash \diamond}
\infer1{\Gamma \vdash \text{continue}: C}
\end{prooftree}
\\
\tag{Command - skip}
\begin{prooftree}
\hypo{\Gamma \vdash \diamond}
\infer1{\Gamma \vdash \text{skip}: C}
\end{prooftree}
\\
\tag{Command - return}
\begin{prooftree}
\hypo{\Gamma \vdash E:A}
\infer1{\Gamma \vdash (\text{return } E): C}
\end{prooftree}
\end{gather}

\subsubsection{Quantification Expressions}
\paragraph{Includes:} Set, sequence, relation, and bag comprehensions, in addition to (set based) lambda expressions. Comprehension variables may be bound through either an $Identifier$ or an arbitrarily nested tuple of $Identifiers$

\begin{gather}
\tag{Lambda Expression}
\begin{prooftree}
\hypo{\Gamma, I_1:T_1, ..., I_n: T_n \vdash E:A, P: bool}
\hypo{I_i \notin dom(\Gamma)}
\hypo{i \neq j \implies I_i \neq I_j}
\infer3{\Gamma \vdash (\lambda I_1, ..., I_n \cdot P \mid E): (T_1, ..., T_n) \pfun A}
\end{prooftree}
\\
\tag{Quantifier}
\begin{prooftree}
\hypo{\Gamma \vdash G: Generator[I]}
\hypo{\Gamma , I \vdash E: A}
\infer2{\Gamma \vdash (G \mid E) : Quantifier[A]}
\end{prooftree}
\\
\tag{Branching Quantifier}
\begin{prooftree}
\hypo{\Gamma \vdash Q_i : Quantifier[A]}
\infer1{\Gamma \vdash (Q_1 ` ... ` Q_n): Quantifier[A]}
\end{prooftree}
\\
\tag{Generator nest}
\begin{prooftree}
\hypo{\Gamma \vdash G_i: Generator[T_i]}
\infer1{\Gamma \vdash G_1, ... G_n : Generator[(T_1,..., T_n)]} % acts like a generator over the cartesian product
\end{prooftree}
\\
\tag{Generator}
\begin{prooftree}
\hypo{\Gamma \vdash S: set[T], P: bool}
\hypo{I = x_1:T_1, ..., x_n:T_n}
\hypo{\Gamma \vdash structuralMatch(I, T)}
\infer3{\Gamma \vdash (I \in S \cdot P): Generator[T]}
\end{prooftree}
\\
% \tag{Quantification Body \footnote{where $TupleIdentifier$ is a recursive tuple of either $Identifier$s or nested $TupleIdentifiers$, $structuralMatch(e,T)$ is a recursive match between identifiers in $e$ and the set type $T$, $hasGenerator(P,e \in g)$ means the expression $e \in g$ occurs in one of the top-level OR/AND clauses, and $structuralMatchInAllORs(e, T, P)$ means all top-level OR clauses must see $e \in g$ for any $g$ with type $T$. For now, we may just force the generator to appear in the first clause of the predicate, where the predicate is a top-level AND.}}
% \begin{prooftree}
% \hypo{I = x_1:T_1, ..., x_n:T_n}
% \hypo{\Gamma \vdash binds(P, I)}
% \hypo{\Gamma, I \vdash P:bool, E:A}
% \infer3{\Gamma \vdash (x_1, ..., x_n \cdot P \mid E): QuantificationBody[I, A]}
% \end{prooftree}
% \\
% \tag{Binds (with generator)}
% \begin{prooftree}
% \hypo{\Gamma \vdash g : set[T]}
% \hypo{\Gamma \vdash structuralMatch(I, T)}
% \infer2{\Gamma \vdash binds(I \in g, I)}
% \end{prooftree}
% \\
% \tag{Binds with OR}
% \begin{prooftree}
% \hypo{\Gamma \vdash binds(P_1, I), ..., binds(P_n, I)}
% \infer1{\Gamma \vdash binds(\bigvee P_i, I)}
% \end{prooftree}
% \\
% \tag{Binds with AND}
% \begin{prooftree}
% \hypo{\Gamma, I \vdash Q: bool}
% \hypo{\Gamma \vdash binds(P, I)}
% \infer2{\Gamma \vdash binds(P \land Q, I)}
% \end{prooftree}
% \\
\tag{Structural Match}
\begin{prooftree}
\hypo{\Gamma \vdash T}
\hypo{x \notin dom(\Gamma)}
\infer2{\Gamma \vdash structuralMatch(x, T)}
\end{prooftree}
\\
\tag{Structural Match with Tuple}
\begin{prooftree}
\hypo{\Gamma \vdash structuralMatch(x_1, T_1), ..., structuralMatch(x_n, T_n)}
\infer1{\Gamma \vdash structuralMatch((x_1, ..., x_n), (T_1, ..., T_n))}
\end{prooftree}
\\
\tag{Set Comprehension}
\begin{prooftree}
\hypo{\Gamma \vdash Q: Quantifier[A]}
\infer1{\Gamma \vdash \setOne{Q}: set[A]}
\end{prooftree}
\\
\tag{Bag Comprehension}
\begin{prooftree}
\hypo{\Gamma \vdash Q: Quantifier[A]}
\infer1{\Gamma \vdash \bagOne{Q}: bag[A]}
\end{prooftree}
\\
\tag{Sequence Comprehension}
\begin{prooftree}
\hypo{\Gamma \vdash Q: Quantifier[A]}
\infer1{\Gamma \vdash \listOne{Q}: sequence[A]}
\end{prooftree}
\\
\tag{Fold}
\begin{prooftree}
\hypo{\Gamma, i: A \vdash Q: Quantifier[A]}
\hypo{\Gamma \vdash Id: A}
\hypo{i \notin dom(\Gamma)}
\infer3{\Gamma \vdash (fold\ i := Id : Q): A}
\end{prooftree}
\\
\tag{Iter Quantifier}
\begin{prooftree}
\hypo{\Gamma \vdash G: IterGenerator[I], x_i:T_i, \text{block }C}
\infer1{\Gamma \vdash (G \rightarrow (x_1, ..., x_n) \mid C): IterQuantifier[(T_1, ..., T_n)]}
\end{prooftree}
\\
\tag{Iter Generator}
\begin{prooftree}
\hypo{\Gamma \vdash G: Generator[I], y_i:T_i}
\hypo{x_i \notin dom(\Gamma)}
\infer2{\Gamma \vdash (G : x_1 := y_1; ... ; x_n := y_n): IterGenerator[I]}
\end{prooftree}
\\
\tag{Iter}
\begin{prooftree}
\hypo{\Gamma \vdash Q: IterQuantifier[A]}
\infer1{\Gamma \vdash (iter\ Q): A}
\end{prooftree}
\end{gather}

Let $\otimes = \{ \bigcup, \bigcap \}$

\begin{gather}
\tag{General Quantifiers}
\begin{prooftree}
\hypo{\Gamma \vdash Q: Quantifier[A]} % where A is a set
\infer1{\Gamma \vdash (\otimes Q): A}
\end{prooftree}
\end{gather}

Let $\otimes = \{ \forall, \exists \}$

\begin{gather}
\tag{Bool Quantifiers}
\begin{prooftree}
\hypo{\Gamma \vdash Q: Quantifier[bool]}
\infer1{\Gamma \vdash (\otimes Q): bool}
\end{prooftree}
\end{gather}

Let $\otimes = \{ \sum, \prod \}$

\begin{gather}
\tag{Numeric Quantifiers}
\begin{prooftree}
\hypo{\Gamma \vdash Q: Quantifier[T]}
\hypo{\Gamma \vdash T \in \{int, nat, float\}}
\infer2{\Gamma \vdash (\otimes Q): T}
\end{prooftree}
\end{gather}

\subsubsection{Boolean Operations}
\paragraph{Includes:} Common predicate logic operations (equivalence, implication, and, or, etc.).

Let $\otimes = \{\equiv, \not\equiv, \implies, \land, \lor \}$

\begin{gather}
\tag{Binary Boolean Operations}
\begin{prooftree}
\hypo{\Gamma \vdash b_1, b_2: bool}
\infer1{\Gamma \vdash b_1 \otimes b_2 : bool}
\end{prooftree}
\\
\tag{Boolean Operations - Negation}
\begin{prooftree}
\hypo{\Gamma \vdash b_1: bool}
\infer1{\Gamma \vdash \lnot b_1 : bool}
\end{prooftree}
\end{gather}

\subsubsection{Equality and Membership}
\paragraph{Includes:} Equality and ordering comparisons, set membership.

Let $\otimes = \{=, \neq\}$

\begin{gather}
\tag{Equals}
\begin{prooftree}
\hypo{\Gamma \vdash a_1: T_1, a_2: T_2}
\infer1{\Gamma \vdash a_1 \otimes a_2 : bool}
\end{prooftree}
\end{gather}

Let $\otimes = \{<, \leq, >, \geq \}$

\begin{gather}
\tag{Ordering Operators}
\begin{prooftree}
\hypo{\Gamma \vdash a_1, a_2: T}
\hypo{T \in \{int, float, nat\}}
\infer2{\Gamma \vdash a_1 \otimes a_2 : bool}
\end{prooftree}
\end{gather}

Let $\otimes = \{\in, \not\in \}$

\begin{gather}
\tag{Set Membership}
\begin{prooftree}
\hypo{\Gamma \vdash a_1: T_1, a_2: set[T_2]}
\infer1{\Gamma \vdash a_1 \otimes a_2: bool}
\end{prooftree}
\end{gather}

\subsubsection{Set Operations}
\paragraph{Includes:} Subset comparisons, union, intersection, cartesian product, etc. Also includes Maplet and Range.

Let $\otimes = \{\subset, \subseteq, \supset, \supseteq, \not\subset, \not\subseteq, \not\supset, \not\supseteq \}$

\begin{gather}
\tag{Set Ordering Operations}
\begin{prooftree}
\hypo{\Gamma \vdash a_1: set[T_1], a_2: set[T_2]}
\infer1{\Gamma \vdash a_1 \otimes a_2: bool}
\end{prooftree}
\end{gather}

Let $\otimes = \{\cup, \cap, \setminus \}$

\begin{gather}
\tag{Set Operations}
\begin{prooftree}
\hypo{\Gamma \vdash a_1, a_2: set[T]}
\infer1{\Gamma \vdash a_1 \otimes a_2: set[T]}
\end{prooftree}
\\
\tag{Cartesian Product}
\begin{prooftree}
\hypo{\Gamma \vdash a_1: set[T_1], a_2: set[T_2]}
\infer1{\Gamma \vdash a_1 \times a_2: T_1 \rel T_2}
\end{prooftree}
\\
\tag{Maplet}
\begin{prooftree}
\hypo{\Gamma \vdash a_1: T_1, a_2: T_2}
\infer1{\Gamma \vdash a_1 \mapsto a_2: tuple[T_1, T_2]}
\end{prooftree}
\\
\tag{Numerical Range}
\begin{prooftree}
\hypo{\Gamma \vdash a_1, a_2: int}
\infer1{\Gamma \vdash a_1 \upto a_2: set[int, \text{ with } min = a_1, \text{ with } max = a_2]}
% Is max/min enough here? what if we union two ranges with dijoint max/mins
\end{prooftree}
\\
\tag{Set Operations - Powerset}
\begin{prooftree}
\hypo{\Gamma \vdash a:T}
\infer1{\Gamma \vdash \mathbb{P}(a):set[T]}
\end{prooftree}
\end{gather}

\subsubsection{Bag Operations}
\paragraph{Includes:} Bag union, bag intersection, etc.

Let $\otimes = \{\cup, \cap, +, -\}$
\begin{gather}
\tag{Bag Operations}
\begin{prooftree}
\hypo{\Gamma \vdash a_1, a_2: bag[T]}
\infer1{\Gamma \vdash a_1 \otimes a_2: bag[T]}
\end{prooftree}
\\
\tag{Bag Operations - Image}
\begin{prooftree}
\hypo{\Gamma \vdash a: T_1 \rel T_2, b: bag[T_1]}
\infer1{\Gamma \vdash a[b]: bag[T_2] \quad a[b] = a^{-1} \circ b}
\end{prooftree}
\end{gather}

\subsubsection{Relation Operations}
\paragraph{Includes:} Relation overriding, composition, domain/range manipulation. See Section 5.1.3 for changes in relational subtypes after applying operations.

\begin{gather}
\tag{Relation Operations - Function Call}
\begin{prooftree}
\hypo{\Gamma \vdash a: T_1 \rel T_2, b:T_1}
\infer1{\Gamma \vdash  a(b):T_2}
\end{prooftree}
\\
\tag{Relation Operations - Image}
\begin{prooftree}
\hypo{\Gamma \vdash a: T_1 \rel T_2, b: set[T_1]}
\infer1{\Gamma \vdash a[b]: set[T_2]}
\end{prooftree}
\\
\tag{Relation Operations - Overriding}
\begin{prooftree}
\hypo{\Gamma \vdash a_1, a_2: T_1 \rel T_2}
\infer1{\Gamma \vdash a_1 \oplus a_2: T_1 \rel T_2}
\end{prooftree}
\\
\tag{Relation Operations - Composition}
\begin{prooftree}
\hypo{\Gamma \vdash a_1: T_1 \rel T_2, a_2: T_2 \rel T_3}
\infer1{\Gamma \vdash a_1 \circ a_2: T_1 \rel T_3}
\end{prooftree}
\\
\tag{Relation Operations - Domain Restriction}
\begin{prooftree}
\hypo{\Gamma \vdash a_1: set[T_1], a_2: T_1 \rel T_2}
\infer1{\Gamma \vdash a_1 \domres a_2: T_1 \rel T_2}
\end{prooftree}
\\
\tag{Relation Operations - Domain Subtraction}
\begin{prooftree}
\hypo{\Gamma \vdash a_1: set[T_1], a_2: T_1 \rel T_2}
\infer1{\Gamma \vdash a_1 \domsub a_2: T_1 \rel T_2}
\end{prooftree}
\\
\tag{Relation Operations - Range Restriction}
\begin{prooftree}
\hypo{\Gamma \vdash a_1: T_1 \rel T_2, a_2: set[T_2]}
\infer1{\Gamma \vdash a_1 \ranres a_2: T_1 \rel T_2}
\end{prooftree}
\\
\tag{Relation Operations - Range Subtraction}
\begin{prooftree}
\hypo{\Gamma \vdash a_1: T_1 \rel T_2, a_2: set[T_2]}
\infer1{\Gamma \vdash a_1 \ransub a_2: T_1 \rel T_2}
\end{prooftree}
\end{gather}

\subsubsection{Relational Subtypes After Operations}
\paragraph{Includes:} Changes in relational subtyping (totality, surjectivity, etc.) after applying operations.
For simplicity, we implement subtype changes in a few binary operators as a boolean mask representing: [with total ($t$), with total on range ($t_r$), with one-to-manyness ($1-$), with many-to-oneness ($1-_r$)]. Subtyping is done conservatively: we keep subtype properties only when they are guaranteed. See Section 5.1.3 for more on relational subtyping.

\begin{gather}
\tag{Relational Subtype - Domain Restriction}
\begin{prooftree}
\hypo{\Gamma \vdash s: set[T_1], r: (T_1 \rel T_2)[True, x, y, z]}
\infer1{\Gamma \vdash (s \domres r): (T_1 \rel T_2)[False, x, y, z]}
\end{prooftree}
\\
\tag{Relational Subtype - Domain Subtraction}
\begin{prooftree}
\hypo{\Gamma \vdash s: set[T_1], r: (T_1 \rel T_2)[True, x, y, z]}
\infer1{\Gamma \vdash (s \domsub r): (T_1 \rel T_2)[False, x, y, z]}
\end{prooftree}
\\
\tag{Relational Subtype - Range Restriction}
\begin{prooftree}
\hypo{\Gamma \vdash s: set[T_2], r: (T_1 \rel T_2)[x, True, y, z]}
\infer1{\Gamma \vdash (r \ranres s): (T_1 \rel T_2)[x, False, y, z]}
\end{prooftree}
\\
\tag{Relational Subtype - Range Subtraction}
\begin{prooftree}
\hypo{\Gamma \vdash s: set[T_2], r: (T_1 \rel T_2)[x, True, y, z]}
\infer1{\Gamma \vdash (r \ransub s): (T_1 \rel T_2)[x, False, y, z]}
\end{prooftree}
\\
\tag{Relational Subtype - Inverse}
\begin{prooftree}
\hypo{\Gamma \vdash r: (T_1 \rel T_2)[w, x, y, z]}
\infer1{\Gamma \vdash r^{-1}: (T_1 \rel T_2)[x, w, z, y]}
\end{prooftree}
\\
\tag{Relational Subtype - Overriding}
\begin{prooftree}
\hypo{\Gamma \vdash r: (T_1 \rel T_2)[a, b, c, d], t: (T_1 \rel T_2)[w, x, y, z]}
\infer1{\Gamma \vdash (r \oplus t): (T_1 \rel T_2)[w, x, c \land y, d \land z]}
\end{prooftree}
\\
\tag{Relational Subtype - Composition}
\begin{prooftree}
\hypo{\Gamma \vdash r: (T_1 \rel T_2)[a, b, c, d], t: (T_2 \rel T_3)[w, x, y, z]}
\infer1{\Gamma \vdash (r \circ t): (T_1 \rel T_3)[a, b \land x, c \land y, d \land z]}
\end{prooftree}
\end{gather}

\subsubsection{Sequence Operations}
\paragraph{Includes:} Concatenation

\begin{gather}
\tag{Sequence Operations - Concatenation}
\begin{prooftree}
\hypo{\Gamma \vdash a_1, a_2: sequence[T]}
\infer1{\Gamma \vdash a_1 \concat a_2: sequence[T]}
\end{prooftree}
\end{gather}

\subsubsection{Numerical Operations}
\paragraph{Includes:} Addition, subtraction, multiplication, etc. on natrual numbers, integers, and floats.

\begin{gather}
\tag{Integer Operations - Division}
\begin{prooftree}
\hypo{\Gamma \vdash a_1: T_1, a_2: T_2}
\hypo{T_i \in \{int, nat\}}
\infer2{\Gamma \vdash a_1 \text{ div } a_2: max(T_1, T_2)}
\end{prooftree}
\\
\tag{Integer Operations - Modulo}
\begin{prooftree}
\hypo{\Gamma \vdash a_1: T_1, a_2: T_2}
\hypo{T_i \in \{int, nat\}}
\infer2{\Gamma \vdash a_1 \text{ mod } a_2: max(T_1, T_2)}
\end{prooftree}
\\
\tag{Numerical Operations - Addition}
\begin{prooftree}
\hypo{\Gamma \vdash a_1: T_1, a_2: T_2}
\hypo{T_i \in \{int, nat, float\}}
\infer2{\Gamma \vdash a_1 + a_2: max(T_1, T_2)}
\end{prooftree}
\\
\tag{Numerical Operations - Subtraction}
\begin{prooftree}
\hypo{\Gamma \vdash a_1: T_1, a_2: T_2}
\hypo{T_i \in \{int, nat, float\}}
\infer2{\Gamma \vdash a_1 - a_2: max(T_1, T_2, int)}
\end{prooftree}
\\
\tag{Numerical Operations - Floating Division}
\begin{prooftree}
\hypo{\Gamma \vdash a_1: T_1, a_2: T_2}
\hypo{T_i \in \{int, nat, float\}}
\infer2{\Gamma \vdash a_1 * a_2: float}
\end{prooftree}
\\
\tag{Numerical Operations - Multiplication}
\begin{prooftree}
\hypo{\Gamma \vdash a_1: T_1, a_2: T_2}
\hypo{T_i \in \{int, nat, float\}}
\infer2{\Gamma \vdash a_1 \text{ div } a_2: max(T_1, T_2)}
\end{prooftree}
\\
\tag{Numerical Operations - Negation}
\begin{prooftree}
\hypo{\Gamma \vdash a_1: T}
\hypo{T \in \{int, nat, float\}}
\infer2{\Gamma \vdash -a_1: max(int, T)}
\end{prooftree}
\\
\tag{Numerical Operations - Exponentiation}
\begin{prooftree}
\hypo{\Gamma \vdash a_1: T_1, a_2: T_2}
\hypo{T_i \in \{int, nat, float\}}
\infer2{\Gamma \vdash a_1 \text{ \^{} } a_2: max(T_1, T_2)}
\end{prooftree}
\end{gather}

\subsubsection{Record Types}
\paragraph{Includes:} Types for named tuples.

\begin{gather}
\tag{Records - Access}
\begin{prooftree}
\hypo{\Gamma \vdash a: Record[I:A, ...]}
\infer1{\Gamma \vdash a.I: A}
\end{prooftree}
\\
\tag{Records - Type Definition}
\begin{prooftree}
\hypo{\Gamma \vdash A_1, ..., A_n}
\hypo{\forall I_i, I_j \cdot I_i \neq I_j}
\infer2{\Gamma \vdash Record[I_1: A_1, ..., I_n: A_n]}
\end{prooftree}
\\
\tag{Records - Initialization}
\begin{prooftree}
\hypo{\Gamma \vdash a_1: A_1, ..., a_n: A_n}
\infer1{\Gamma \vdash record(a_1=A_1, ..., a_n=A_n): Record[a_1: A_1, ..., a_n:A_n]}
\end{prooftree}
\end{gather}

\subsubsection{Compound Commands}
\paragraph{Includes:} If, While, For, and Import statements, along with Procedure definitions and calls.

\begin{gather}
\tag{Command - Composition}
\begin{prooftree}
\hypo{\Gamma \vdash C_1, C_2}
\infer1{\Gamma \vdash C_1 \text{ \textbackslash n } C_2}
\end{prooftree}
\\
\tag{Command - Valid Import Module}
\begin{prooftree}
\hypo{\Gamma \vdash \diamond}
\infer1{\Gamma \vdash \text{import } M: C}
\end{prooftree}
\\
\tag{Command - Valid Import Names}
\begin{prooftree}
\hypo{\Gamma \vdash \diamond}
\infer1{\Gamma \vdash \text{from } M \text{ import } I_1, ..., I_n: C}
\end{prooftree}
\\
\tag{Command - Import Module}
\begin{prooftree}
\hypo{\Gamma \vdash \text{import } M \text{ (with identifiers and types } I_1: A_1, ..., I_n: A_n)}
\infer1{\Gamma, M=record(I_1=A_1, ..., I_n=A_n) \vdash \diamond}
\end{prooftree}
\\
\tag{Command - Import Names}
\begin{prooftree}
\hypo{\Gamma \vdash \text{from } M \text{ import } I_1, ..., I_n \text{ (with types } A_1, ..., A_n)}
\infer1{\Gamma, I_1:A_1, ..., I_n:A_n \vdash \diamond}
\end{prooftree}
\\
\tag{Command - Procedure Definition}
\begin{prooftree}
% Read this as: make a new scope, place the procedure I and params x_* to create the body C:R.
% Then, the current scope produces a symbol I with type procedure[...]
\hypo{\Gamma, I: procedure[A_1, ..., A_n, R], x_1: A_1, ..., x_n: A_n \vdash C: R}
\hypo{I, x_1, ..., x_n \notin dom(\Gamma)}
\hypo{i \neq j \implies x_i \neq x_j}
\infer3{\Gamma \vdash \text{ procedure }I=C : procedure[A_1, ..., A_n, R]}
\end{prooftree}
\\
\tag{Command - Procedure Call}
\begin{prooftree}
\hypo{\Gamma \vdash procedure[A_1, ..., A_n, R], a_1: A_1, ..., a_n:A_n}
\infer1{\Gamma \vdash I(a_1, ..., a_n): R}
\end{prooftree}
\\
\tag{Command - If}
\begin{prooftree}
% C is a new scope (block) but doesnt necessarily define new bindings
\hypo{\Gamma \vdash b: bool, \text{block }C}
\infer1{\Gamma \vdash \text{if } b: C}
\end{prooftree}
\\
\tag{Command - If Else}
\begin{prooftree}
\hypo{\Gamma \vdash b: bool, \text{block }C_1, \text{block }C_2}
\infer1{\Gamma \vdash \text{if } b: C_1 \text{ else}: C_2}
\end{prooftree}
\\
\tag{Command - While}
\begin{prooftree}
\hypo{\Gamma \vdash b: bool, \text{block }C}
\infer1{\Gamma \vdash \text{while } b: C}
\end{prooftree}
\\
\tag{Command - For}
\begin{prooftree}
\hypo{\Gamma \vdash G: set[T]}
\hypo{\Gamma, i: T \vdash C}
\hypo{i \notin dom(\Gamma)}
\infer3{\Gamma \vdash \text{for } i \in G: C}
\end{prooftree}
\\
\tag{Command - Block}
\begin{prooftree}
% Declarations D produce a new environment in C
\hypo{\Gamma \vdash D \therefore \Gamma'}
\hypo{\Gamma' \vdash C}
\infer2{\Gamma \vdash \text{block }D \text{ \textbackslash n } C}
\end{prooftree}
\end{gather}



\subsubsection{Built-in Functions}
\paragraph{Includes:} Common set-related functions.

\begin{gather}
\tag{Built-in - Minimum}
\begin{prooftree}
\hypo{\Gamma \vdash T[\text{with } order \in traits(T)]}
\infer1{\Gamma \vdash min: set[T] \rightarrow T}
\end{prooftree}
\\
\tag{Built-in - Mapped Minimum}
\begin{prooftree}
\hypo{\Gamma \vdash T, int}
\infer1{\Gamma \vdash min: set[T] \times (T \rightarrow int) \rightarrow T}
\end{prooftree}
\\
\tag{Built-in - Maximum}
\begin{prooftree}
\hypo{\Gamma \vdash T[\text{with } order \in traits(T)]}
\infer1{\Gamma \vdash max: set[T] \rightarrow T}
\end{prooftree}
\\
\tag{Built-in - Mapped Maximum}
\begin{prooftree}
\hypo{\Gamma \vdash T, int}
\infer1{\Gamma \vdash max: set[T] \times (T \rightarrow int) \rightarrow T}
\end{prooftree}
\\
\tag{Built-in - Choice}
\begin{prooftree}
\hypo{\Gamma \vdash T}
\infer1{\Gamma \vdash choice: set[T] \rightarrow T}
\end{prooftree}
\\
\tag{Built-in - Domain}
\begin{prooftree}
\hypo{\Gamma \vdash T}
\infer1{\Gamma \vdash dom: T_1 \rel T_2 \rightarrow set[T_1]}
\end{prooftree}
\\
\tag{Built-in - Range}
\begin{prooftree}
\hypo{\Gamma \vdash T}
\infer1{\Gamma \vdash ran: T_1 \rel T_2 \rightarrow set[T_2]}
\end{prooftree}
\\
\tag{Built-in - Cardinality}
\begin{prooftree}
\hypo{\Gamma \vdash T}
\infer1{\Gamma \vdash card: set[T] \rightarrow nat}
\end{prooftree}
\\
\tag{Built-in - Bag Size}
\begin{prooftree}
\hypo{\Gamma \vdash T}
\infer1{\Gamma \vdash size: bag[T] \rightarrow nat}
\end{prooftree}
\\
\tag{Built-in - Sum}
\begin{prooftree}
\hypo{\Gamma \vdash T}
\hypo{T \in \{int, float, nat\}}
\infer2{\Gamma \vdash sum: set[T] \rightarrow T}
\end{prooftree}
\\
\tag{Built-in - Cast}
\begin{prooftree}
\hypo{\Gamma \vdash T_1, T_2}
\infer1{\Gamma \vdash cast: type[T_1] \times T_1 \rightarrow T_2}
\end{prooftree}
\\
\tag{Built-in - Cast With \footnote{$E$ represents `with' expressions to tag onto a type}}
\begin{prooftree}
\hypo{\Gamma \vdash T_1, T_2, x: E}
\infer1{\Gamma \vdash cast: type[T_1] \times T_1 \times set[E] \rightarrow T_2[\text{with } x: E]}
\end{prooftree}
\end{gather}

\subsection{Static Analysis - Type System Extensions with Traits}

% What are traits?
Traits extend Simile's type system by giving programmers the ability to provide additional context for the intended use of a defined variable. Traits never override Simile's set theoretic semantics; a program with no user-defined traits is assumed to be correct and functional. However, traits can influence the optimization process and concretization of abstract data structures.  With user-provided context, static analysis becomes more complete and the optimizer may choose more aggressive rewrites. In a sense, traits in Simile are akin to pragmas in other languages, but closer to type refinements rather than direct compiler commands.

For example, consider a variable defining a relation:
\begin{algorithmic}
  \State $r: relation[int, int] := \{\}$
\end{algorithmic}
Without insight into the use of the variable $r$, we may be forced to concretize the abstract type conservatively. The compiler may choose the safest option: straightforward array of tuples or a bidirectional hashmap, with both concrete structures giving up time complexity or memory efficiency respectively. However, if we as domain experts know that this relation is one-to-one, we can provide such insights to the compiler directly:
\begin{algorithmic}
  \State $r: relation[int, int] := \{\}$
  \State \quad trait $many\_to\_one$
  \State \quad trait $one\_to\_many$
\end{algorithmic}
Concretization can then see that a hashmap is well-suited to representing the data. Even further, since $r$ is one-to-one, we may even choose to save memory by concretizing to two arrays aligned by index. To the programmer, the important semantics remain the same; $r$ is a relation and behaves like one for the duration of the program, but the compiler can choose to swap the detailed implementation with one better-suited for the problem.

There is one caveat to building these features into the language itself: we trust that users will not provide deliberately misleading traits. It is difficult to check that trait assignments themselves are correct or well-defined through static analysis alone, so an unknowing programmer may stumble into undefined behaviour. For example, consider:
\begin{algorithmic}
  \State $x: int := 3$
  \State \quad trait $minimum = 1$
  \State \quad trait $maximum = 4$
\end{algorithmic}
The compiler may see that this integer is restricted to a very small range and thus choose a smaller integer representation size. However, if the programmer then writes:
\begin{algorithmic}
  \State $x := 5$
\end{algorithmic}
This assumption breaks and does not fit into the general type. We may be able to spot such hardcoded misuses of the trait system, but if $5$ was instead an I/O function, we can only provide runtime checks that slow down the system (and may turn out to be slower than just using a normal i32 type in the first place.

Given the power and danger of using traits, programs without traits by default should behave as expected. Note that this danger happens only on assignment, since traits are assumed correct for the variable to which they are assigned. When performing operations, the compiler will widen or narrow traits as is needed to remain in a good state.

\subsubsection{Available Traits}
The following paragraphs list and categorize identified traits.

For parameterized types that accept type arguments, traits will be propagated to the ``parent'' type as much as possible. For example, a $set[int]$ restricting the range of the underlying integers should have that trait specified on the outer $set$ type rather than the underlying $int$ type.
% For parameterized types that accept type arguments (like set, relation, sequence, tuple, ...), some traits may apply to the type argument when sensible. For example, a $set[int]$ restricting the range of the underlying integers should have that trait specified on the underlying $int$ type rather than the outer $set$ type (although we may loosen this with syntactic sugar). Like Event-B, $set$ types are really powersets of the underlying element type, whereas primitive types represent values chosen from a \textit{set} of primitive values.

\paragraph{Base Traits} Applicable to almost every object
\begin{itemize}
    \item \textit{Literal}: A trait containing a single literal value. If a type has this trait, it is guaranteed to be of this value.
    \item \textit{Domain}: A looser form of \textit{Literal} that carries a set of possible values. When applied, the type is restricted to (one of) the values from this $Domain$. Relations use a different definition of domain.
    \item \textit{Undefined}: Intended for use by compiler analysis. Used when a value \textit{may} be undefined. Ex. $choice(\{\})$.
    % \item Leave out for now: $UnknownElementInDomain$: Loosens the $Domain$ guarantee so that values in the domain are only likely to appear, but not guaranteed? We would need to see if this opens up any optimizations...
    \item \textit{Immutable}: Represents values that cannot be reassigned/modified. %Should we instead assume immutability and have a Mutable trait?
\end{itemize}

\paragraph{Arithmetic} Restricts the value of primitive types
\begin{itemize}
    \item \textit{Orderable}: Values of this type can be ordered. This may be useful for order-based implementations like BTrees for sets.
    \item \textit{Min}: Minimum value.
    \item \textit{Max}: Maximum value. Combine with the \textit{Min} trait to create a range.
\end{itemize}

\paragraph{Collection Traits}
\begin{itemize}
    \item \textit{Iterable}: All collections are iterable (they can be used as a for-loop generator).
    \item \textit{Empty}: A collection with guaranteed cardinality = 0.
    \item \textit{Size}: An over-estimate of a set's size. When available, Size will never under-estimate the number of elements in a collection.
    \item \textit{Unique}: All entries are guaranteed to be unique. More useful for non-set collections (which are always assumed to be unique). Ex. a non-repeating sequence.
    \item \textit{Total}: The elements in a set cover the entirety of the domain of the underlying element type. Relations have their own definition of totality.
\end{itemize}

\paragraph{Relations}
Represents the relational subtype properties (injectiveness, surjectiveness, etc.)
\begin{itemize}
    \item \textit{ManyToOne}: Elements of the domain may be repeated with different range values.
    \item \textit{OneToMany}: Elements of the range may be repeated with different range values. \textit{ManyToOne} and \textit{OneToMany} imply a one-to-one mapping.
    \item \textit{TotalOnDomain}: All elements of the domain are within the mapping.
    \item \textit{TotalOnRange}: All elements of the range are within the mapping.
    \item \textit{RelationalDomain}: The set of elements that may populate the relation's domain. Replacement of \textit{Domain} for relations.
    \item \textit{RelationalRange}: The set of elements that may populate the relation's range.
\end{itemize}

\paragraph{Generics}
\begin{itemize}
    \item \textit{GenericBound}: Contains the set of types a generic may become. For example, the following generic type represents any numeric value:
    \begin{algorithmic}
      \State $T: type := generic$
      \State \quad trait $generic\_bound = int$
      \State \quad trait $generic\_bound = float$
    \end{algorithmic}
\end{itemize}

\paragraph{Procedure}
\begin{itemize}
    \item \textit{TreatAsExpr}: Inline this function and treat it as an expression/operator. Useful for having function-call notation act like an operator. For example:
    \begin{algorithmic}
      \State procedure $add(s: set[T], e: T) -> set[T]:$
      \State \quad trait $treat\_as\_expr$
      \State \quad return $s \cup \{e\}$
    \end{algorithmic}
\end{itemize}

\subsubsection{Type-trait interactions}
Some types are incompatible with some traits, and some types must have some traits. We provide tables showing compatibility between traits and types. A double-checkmark means the trait is always present - a base trait of that type.

\paragraph{Primitives}
\begin{center}
\begin{longtable}{l c c c c c}
    & none & bool & int & float & str \\
    \endhead
    Immutable & \cmark & \cmark & \cmark & \cmark & \cmark \\
    Literal  & None & \cmark & \cmark & \cmark & \cmark \\
    Undefined  & \cmark & \cmark & \cmark & \cmark & \cmark \\
    Orderable  & \xmark & \xmark & \ccmark & \ccmark & \ccmark \\ % str uses lexicographic ordering
    Min  & \xmark & \xmark & \cmark & \cmark & \cmark \\
    Max  & \xmark & \xmark & \cmark & \cmark & \cmark \\
    Domain  & \{None\} & \{True, False\} & \cmark & \cmark & \cmark \\
    Iterable  & \xmark & \xmark & \xmark & \xmark & \ccmark \\
    Unique  & \xmark & \xmark & \xmark & \xmark & \xmark \\
    Empty  & \xmark & \xmark & \xmark & \xmark & \cmark \\
    Size  & \xmark & \xmark & \xmark & \xmark & \cmark \\
    Total  & \xmark & \xmark & \xmark & \xmark & \xmark \\
    OneToMany  & \xmark & \xmark & \xmark & \xmark & \xmark \\
    ManyToOne  & \xmark & \xmark & \xmark & \xmark & \xmark \\
    TotalOnDomain  & \xmark & \xmark & \xmark & \xmark & \xmark \\
    TotalOnRange  & \xmark & \xmark & \xmark & \xmark & \xmark \\
    RelationalDomain  & \xmark & \xmark & \xmark & \xmark & \xmark \\
    RelationalRange  & \xmark & \xmark & \xmark & \xmark & \xmark \\
    GenericBound  & \xmark & \xmark & \xmark & \xmark & \xmark \\
    TreatAsExpr  & \xmark & \xmark & \xmark & \xmark & \xmark \\
\end{longtable}
\end{center}

\paragraph{Collections}
\begin{center}
\begin{longtable}{l c c c c c c}
    & tuple & set & enum & relation & bag & sequence \\
    \endhead
    Immutable & \cmark & \cmark & \ccmark & \cmark & \cmark & \cmark \\ % TODO should add frozen to restrict modification of the set itself, rather than just restricting reassignment
    Literal & \cmark & \cmark & \cmark & \cmark & \cmark & \cmark \\
    Undefined & \cmark & \cmark & \cmark & \cmark & \cmark & \cmark \\
    Orderable & \cmark & \cmark & \cmark & \cmark & \cmark & \cmark \\ % For most of these, only orderable if the underlying types are orderable
    Min & \cmark & \cmark & \cmark & \cmark & \cmark & \cmark \\
    Max & \cmark & \cmark & \cmark & \cmark & \cmark & \cmark \\
    Domain & \cmark & \cmark & \ccmark & \xmark & \cmark & \cmark \\
    Iterable & \ccmark & \ccmark & \ccmark & \ccmark & \ccmark & \ccmark \\
    Unique & \cmark & \ccmark & \ccmark & \cmark & \cmark & \cmark \\
    Empty & \cmark & \cmark & \cmark & \cmark & \cmark & \cmark \\
    Size & \cmark & \cmark & \ccmark & \cmark & \cmark & \cmark \\ % TODO add a count for bags (the sum of range values)
    Total & \cmark & \cmark & \cmark & \xmark & \cmark & \cmark \\
    OneToMany & \xmark & \xmark & \xmark & \cmark & \cmark & \cmark \\
    ManyToOne & \xmark & \xmark & \xmark & \cmark & \ccmark & \ccmark \\
    TotalOnDomain & \xmark & \xmark & \xmark & \cmark & \xmark & \xmark \\
    TotalOnRange & \xmark & \xmark & \xmark & \cmark & \xmark & \xmark \\
    RelationalDomain & \xmark & \xmark & \xmark & \cmark & \xmark & \xmark \\
    RelationalRange & \xmark & \xmark & \xmark & \cmark & \xmark & \xmark \\
    GenericBound & \xmark & \xmark & \xmark & \xmark & \xmark & \xmark \\
    TreatAsExpr & \xmark & \xmark & \xmark & \xmark & \xmark & \xmark \\
\end{longtable}
\end{center}

\paragraph{Other}
\begin{center}
\begin{longtable}{l c c c}
    & record & procedure & type (generic) \\
    \endhead
    Immutable & \cmark & \xmark & \cmark \\
    Literal & \cmark & \xmark & \cmark \\
    Undefined & \cmark & \cmark & \cmark \\
    Orderable & \xmark & \xmark & \cmark \\
    Min & \xmark & \xmark & \cmark \\
    Max & \xmark & \xmark & \cmark \\
    Domain & \xmark & \xmark & \cmark \\ %
    % How should procedure handle domains? similar to relations? perhaps it can use its own tuple domain, like trait arg_domain[0] = {...};trait arg_domain[1] = {...};
    % or we could do arg_domain = ({...}, {...}, {...}, ...)
    Iterable & \ccmark & \xmark & \cmark \\
    Unique & \xmark & \xmark & \cmark \\
    Empty & \xmark & \xmark & \cmark \\
    Size & \ccmark & \xmark & \cmark \\
    Total & \xmark & \xmark & \cmark \\
    OneToMany & \xmark & \xmark & \cmark \\
    ManyToOne & \ccmark & \xmark & \cmark \\
    TotalOnDomain & \ccmark & \xmark & \cmark \\
    TotalOnRange & \xmark & \xmark & \cmark \\ % Records are not guaranteed to have the same "range" type across all fields
    RelationalDomain & \ccmark (field names) & \xmark & \cmark \\
    RelationalRange & \xmark & \xmark & \cmark \\
    GenericBound & \xmark & \xmark & \cmark \\
    TreatAsExpr & \xmark & \cmark & \cmark \\ % Generics can be bound to procedures too...
\end{longtable}
\end{center}

\subsubsection{Trait-trait Interactions}
Traits can interact with other traits outside of types. Universal interactions, which don't follow a type or its operations, generally fall into two categories:
\begin{itemize}
    \item Derivable traits. Often, the jurisdiction of traits will overlap with one another: a Literal trait implies a \textit{Domain}, and we can derive the \textit{Min}/\textit{Max} traits from an \textit{Orderable} \textit{Domain}. These will be automatically derived by the compiler when needed.
    \item Incompatible traits. Interactions between these traits are unresolvable automatically and would leave the trait system in a bad state if left alone. To handle incompatible traits safely, we refuse to compile the program until the issue is resolved. Automatically derived traits should never become incompatible (only user-defined traits clashing with user/automatic traits).
\end{itemize}

Most traits are simple flags, as in, ``does this type have this property?''. In the following type (or rather, trait) rules, trait names used with boolean operations should be interpreted as truthy (it is present). A type can only have one of each trait to prevent ambiguity or clashing.

% TODO add in semantics for multiple trait definitions on one assignment
% TODO generic_bound should only be definable once and take in a set of types? dont allow repeated definitions

% OLD TODO design decision: should traits be applied automatically or just verified? If we apply automatically, we could give best effort and then error/remove traits that are incompatible. In verification, we just throw an error. We may use trait restrictions per-type to further determine incompatibilities (ex. an integer cannot be empty?).
\paragraph{Derivable Traits}
% TODO fix below according to code

\begin{gather}
  % \tag{Literal Empty}
  % \begin{prooftree}
  %   \hypo{Literals = \{\}}
  %   \infer1{\lnot Literals}
  % \end{prooftree}
  % \\
  \tag{Domain Empty}
  \begin{prooftree}
    \hypo{Domain = \{\}}
    \infer1{\lnot Domain}
  \end{prooftree}
  \\
  \tag{Literal implies a Domain (under primitive)}
  \begin{prooftree}
    \hypo{Literal}
    \hypo{Literal.type \in \{bool, int, float, string, none, tuple\}}
    \hypo{\lnot Domain}
    \infer3{Domain := \{Literal\}}
  \end{prooftree}
  \\
  \tag{Literal implies a Domain (under collection)}
  \begin{prooftree}
    \hypo{Literal}
    \hypo{Literal.type \in \{set, relation, bag, sequence\}}
    \hypo{\lnot Domain}
    \infer3{Domain := Literal}
  \end{prooftree}
  \\
  \tag{Literal within Domain}
  \begin{prooftree}
    \hypo{Literal}
    \hypo{Domain}
    \infer2{Domain := Domain \cup {Literal}}
  \end{prooftree}
  \\
  \tag{Orderable Domain without Min}
  \begin{prooftree}
    \hypo{Domain}
    \hypo{Orderable}
    \hypo{\lnot Min}
    \infer3{Min := min(Domain)}
  \end{prooftree}
  \\
  \tag{Orderable Domain with Min}
  \begin{prooftree}
    \hypo{Domain}
    \hypo{Orderable}
    \hypo{Min}
    \infer3{Min := min(Min, Domain)}
  \end{prooftree}
  \\
  \tag{Orderable Domain without Max}
  \begin{prooftree}
    \hypo{Domain}
    \hypo{Orderable}
    \hypo{\lnot Max}
    \infer3{Max = max(Domain)}
  \end{prooftree}
  \\
  \tag{Orderable Domain with Max}
  \begin{prooftree}
    \hypo{Domain}
    \hypo{Orderable}
    \hypo{Max}
    \infer3{Max = max(Max, Domain)}
  \end{prooftree}
  \\
  \tag{Min implies Order}
  \begin{prooftree}
    \hypo{Min}
    \infer1{Orderable}
  \end{prooftree}
  \\
  \tag{Max implies Order}
  \begin{prooftree}
    \hypo{Max}
    \infer1{Orderable}
  \end{prooftree}
  \\
  \tag{Full set}
  \begin{prooftree}
    \hypo{Unique}
    \hypo{Size = len(Domain)}
    \hypo{\lnot Empty} % could an empty set be total if its domain consists of nothing? what is this a set of?
    \infer3{Total}
  \end{prooftree}
  \\
  \tag{Empty Size}
  \begin{prooftree}
    \hypo{Size = 0}
    \infer1{Empty}
  \end{prooftree}
  \\
  \tag{Non-Empty Size}
  \begin{prooftree}
    \hypo{Size \neq 0}
    \infer1{\lnot Empty}
  \end{prooftree}
  \\
  \tag{Size implies Iterable}
  \begin{prooftree}
    \hypo{Size}
    \infer1{Iterable}
  \end{prooftree}
  \\
  \tag{Orderable Literal is Min}
  \begin{prooftree}
    \hypo{Literal}
    \hypo{Orderable}
    \hypo{\lnot Min}
    \infer3{Min=Literal}
  \end{prooftree}
  \\
  \tag{Orderable Literal is Max}
  \begin{prooftree}
    \hypo{Literal}
    \hypo{Orderable}
    \hypo{\lnot Max}
    \infer3{Max=Literal}
  \end{prooftree}
  % \\
  % \tag{Unknown value clears literals TODO move this to the operation section TODO add unknown trait so we can keep domain + unknown? Could help some values...}
  % \begin{prooftree}
  %   \hypo{LiteralTraits}
  %   \hypo{unknown\_operation}
  %   \infer2{\lnot LiteralTraits}
  % \end{prooftree}
\end{gather}

\paragraph{Incompatible Traits}
% TODO

\subsubsection{Operation-Trait Interactions}
Operations on \textit{Literal} traits will always produce a \textit{Literal} where possible, following the effects of applying the function on values. For conciseness, we omit all \textit{Literal} interactions and instead refer to the semantics of the operations themselves. Further, we use a shorthand for defining traits on types: $T[<trait>]$ is the same as saying:
\begin{algorithmic}
    \State $T: type := generic$
    \State \quad trait $<trait>$
\end{algorithmic}
% Which has yet to be written formally...

General notes:
\begin{itemize}
  \item Any operation with an undefined value will have an undefined result.
  \item Any operations with immutable inputs will have an immutable output. If multi-arg operations have at least one mutable input, the output will also be mutable (unless immutability appears in the traits below).
\end{itemize}

\paragraph{Base}
\subparagraph{Equals} Not equals is just not applied to equals.
\begin{gather}
  \tag{Literal}
  \begin{prooftree}
    \hypo{a: X[literal=x]}
    \hypo{b: Y[literal=y]}
    \infer2{(a = b): Bool[literal=(X=Y \land x=y)]}
  \end{prooftree}
\end{gather}
\paragraph{Bool}
\subparagraph{Not}
\begin{gather}
  \tag{Literal}
  \begin{prooftree}
    \hypo{p: Bool[literal=x]}
    \infer1{(\lnot p): Bool[literal=(\lnot x)]}
  \end{prooftree}
\end{gather}
\subparagraph{Equivalent} Not equivalent is just not applied to the equivalence operation.
\begin{gather}
  \tag{Literal}
  \begin{prooftree}
    \hypo{a: Bool[literal=x]}
    \hypo{b: Bool[literal=y]}
    \infer2{(a \equiv b): Bool[literal=(x=y)]}
  \end{prooftree}
\end{gather}
\subparagraph{Implies}
\begin{gather}
  \tag{Literal}
  \begin{prooftree}
    \hypo{a: Bool[literal=x]}
    \hypo{b: Bool[literal=y]}
    \infer2{(a \implies b): Bool[literal=(x \implies y)]}
  \end{prooftree}
\end{gather}
\subparagraph{And}
\begin{gather}
  \tag{Literal}
  \begin{prooftree}
    \hypo{a: Bool[literal=x]}
    \hypo{b: Bool[literal=y]}
    \infer2{(a \land b): Bool[literal=(x \land y)]}
  \end{prooftree}
  \\
  \tag{Undefined with false}
  \begin{prooftree}
    \hypo{a: Bool[undefined]}
    \hypo{b: Bool[literal=false]}
    \infer2{(a \land b): Bool[literal=false]}
  \end{prooftree}
  \\
  \tag{Undefined with true}
  \begin{prooftree}
    \hypo{a: Bool[undefined]}
    \hypo{b: Bool[literal=true]}
    \infer2{(a \land b): Bool[undefined]}
  \end{prooftree}
  \\
  \tag{Undefined is not guaranteed}
  \begin{prooftree}
    \hypo{a: Bool[undefined]}
    \hypo{b: Bool}
    \infer2{(a \land b): Bool}
  \end{prooftree}
\end{gather}
\subparagraph{Or}
\begin{gather}
  \tag{Literal}
  \begin{prooftree}
    \hypo{a: Bool[literal=x]}
    \hypo{b: Bool[literal=y]}
    \infer2{(a \lor b): Bool[literal=(x \lor y)]}
  \end{prooftree}
  \\
  \tag{Undefined with false}
  \begin{prooftree}
    \hypo{a: Bool[undefined]}
    \hypo{b: Bool[literal=false]}
    \infer2{(a \lor b): Bool[undefined]}
  \end{prooftree}
  \\
  \tag{Undefined with true}
  \begin{prooftree}
    \hypo{a: Bool[undefined]}
    \hypo{b: Bool[literal=true]}
    \infer2{(a \lor b): Bool[literal=true]}
  \end{prooftree}
  \\
  \tag{Undefined is not guaranteed}
  \begin{prooftree}
    \hypo{a: Bool[undefined]}
    \hypo{b: Bool}
    \infer2{(a \lor b): Bool}
  \end{prooftree}
\end{gather}
\paragraph{Arithmetic}
\subparagraph{Less than} Opposite for greater than.
\begin{gather}
  \tag{Literal}
  \begin{prooftree}
    \hypo{a: T[literal=x, orderable]}
    \hypo{b: T[literal=y, orderable]}
    \hypo{x < y}
    \infer3{(a < b): Bool[literal=(x < y)]}
  \end{prooftree}
  \\
  \tag{Max-min literal true}
  \begin{prooftree}
    \hypo{a: T[max=x]}
    \hypo{b: T[min=y]}
    \hypo{x < y}
    \infer3{(a < b): Bool[literal=true]}
  \end{prooftree}
  \\
  \tag{Max-min literal false}
  \begin{prooftree}
    \hypo{a: T[min=x]}
    \hypo{b: T[max=y]}
    \hypo{x \geq y}
    \infer3{(a < b): Bool[literal=false]}
  \end{prooftree}
\end{gather}
\subparagraph{Less than or equals} Opposite for greater than or equals.
\begin{gather}
  \tag{Literal}
  \begin{prooftree}
    \hypo{a: T[literal=x, orderable]}
    \hypo{b: T[literal=y, orderable]}
    \hypo{x \leq y}
    \infer3{(a \leq b): Bool[literal=(x \leq y)]}
  \end{prooftree}
  \\
  \tag{Less than or equal with max-min}
  \begin{prooftree}
    \hypo{a: T[max=x]}
    \hypo{b: T[min=y]}
    \hypo{x \leq y}
    \infer3{(a \leq b): Bool[literal=true]}
  \end{prooftree}
  \\
  \tag{Less than or equal with min-max}
  \begin{prooftree}
    \hypo{a: T[min=x]}
    \hypo{b: T[max=y]}
    \hypo{x > y}
    \infer3{(a \leq b): Bool[literal=false]}
  \end{prooftree}
\end{gather}
\subparagraph{Negate}
\begin{gather}
  \tag{Literal}
  \begin{prooftree}
    \hypo{p: T[literal=x]}
    \hypo{T \in \{Int, Float\}}
    \infer2{(- p): T[literal=(- x)]}
  \end{prooftree}
  \\
  \tag{Min}
  \begin{prooftree}
    \hypo{p: T[min=x]}
    \hypo{T \in \{Int, Float\}}
    \infer2{(- p): T[max=(- x)]}
  \end{prooftree}
  \\
  \tag{Max}
  \begin{prooftree}
    \hypo{p: T[max=x]}
    \hypo{T \in \{Int, Float\}}
    \infer2{(- p): T[min=(- x)]}
  \end{prooftree}
  \\
  \tag{Domain}
  \begin{prooftree}
    \hypo{p: T[domain=s]}
    \hypo{T \in \{Int, Float\}}
    \infer2{(- p): T[domain=\bSetT{x in s}{-x}]}
  \end{prooftree}
\end{gather}
\subparagraph{Int divide}
\begin{gather}
  \tag{Literal}
  \begin{prooftree}
    \hypo{a: Int[literal=x]}
    \hypo{b: Int[literal=y]}
    \infer2{(a \text{ div } b): T[literal=(\lfloor x / y \rfloor)]}
  \end{prooftree}
  \\
  \tag{Literal zero}
  \begin{prooftree}
    \hypo{a: Int[literal=0]}
    \hypo{b: Int[0 \not\in domain \lor max < 0 \lor min > 0]}
    \infer2{(a \text{ div } b): T[literal=0]}
  \end{prooftree}
  \\
  \tag{Min and literal}
  \begin{prooftree}
    \hypo{a: Int[min=x]}
    \hypo{b: Int[literal=y]}
    \hypo{x \geq 0}
    \infer3{(a \text{ div } b): T[min=(\lfloor x / y \rfloor)]}
  \end{prooftree}
  \\
  \tag{Max and literal}
  \begin{prooftree}
    \hypo{a: Int[max=x]}
    \hypo{b: Int[literal=y]}
    \hypo{x \geq 0}
    \infer3{(a \text{ div } b): T[max=(\lfloor x / y \rfloor)]}
  \end{prooftree}
  \\
  \tag{Literal and min}
  \begin{prooftree}
    \hypo{a: Int[literal=x]}
    \hypo{b: Int[min=y]}
    \hypo{x > 0}
    \hypo{y > 0}
    \infer4{(a \text{ div } b): T[max=(\lfloor x / y \rfloor)]}
  \end{prooftree}
  \\
  \tag{Literal and max}
  \begin{prooftree}
    \hypo{a: Int[literal=x]}
    \hypo{b: Int[max=y]}
    \hypo{x > 0}
    \hypo{y < 0}
    \infer4{(a \text{ div } b): T[min=(\lfloor x / y \rfloor)]}
  \end{prooftree}
  \\
  \tag{Negative literal and min}
  \begin{prooftree}
    \hypo{a: Int[literal=x]}
    \hypo{b: Int[min=y]}
    \hypo{x < 0}
    \hypo{y > 0}
    \infer4{(a \text{ div } b): T[min=(\lfloor x / y \rfloor)]}
  \end{prooftree}
  \\
  \tag{Negative literal and max}
  \begin{prooftree}
    \hypo{a: Int[literal=x]}
    \hypo{b: Int[max=y]}
    \hypo{x < 0}
    \hypo{y < 0}
    \infer4{(a \text{ div } b): T[max=(\lfloor x / y \rfloor)]}
  \end{prooftree}
\end{gather}
\subparagraph{Modulo}
\begin{gather}
  \tag{Literal}
  \begin{prooftree}
    \hypo{a: Int[literal=x]}
    \hypo{b: Int[literal=y]}
    \infer2{(a \text{ mod } b): T[literal=(x \text{ mod } y)]}
  \end{prooftree}
  \\
  \tag{Positive literal}
  \begin{prooftree}
    \hypo{a: Int}
    \hypo{b: Int[max=x]}
    \hypo{x > 0}
    \infer3{(a \text{ mod } b): T[max=x, min=0]}
  \end{prooftree}
  \\
  \tag{Negative literal}
  \begin{prooftree}
    \hypo{a: Int}
    \hypo{b: Int[min=x]}
    \hypo{x < 0}
    \infer3{(a \text{ mod } b): T[min=x, max=0]}
  \end{prooftree}
\end{gather}
\subparagraph{Add} And commutative versions.
\begin{gather}
  \tag{Literal}
  \begin{prooftree}
    \hypo{a: T[literal=x]}
    \hypo{b: T[literal=y]}
    \hypo{T \in \{Int, Float\}}
    \infer3{(a + b): T[literal=(x + y)]}
  \end{prooftree}
  \\
  \tag{Literal and min}
  \begin{prooftree}
    \hypo{a: T[literal=x]}
    \hypo{b: T[min=y]}
    \hypo{T \in \{Int, Float\}}
    \infer3{(a + b): T[min=(x + y)]}
  \end{prooftree}
  \\
  \tag{Literal and max}
  \begin{prooftree}
    \hypo{a: T[literal=x]}
    \hypo{b: T[max=y]}
    \hypo{T \in \{Int, Float\}}
    \infer3{(a + b): T[max=(x + y)]}
  \end{prooftree}
  \\
  \tag{Min}
  \begin{prooftree}
    \hypo{a: T}
    \hypo{b: T[min=y]}
    \hypo{T \in \{Int, Float\}}
    \infer3{(a + b): T[min=y]}
  \end{prooftree}
  \\
  \tag{Max}
  \begin{prooftree}
    \hypo{a: T}
    \hypo{b: T[max=y]}
    \hypo{T \in \{Int, Float\}}
    \infer3{(a + b): T[max=y]}
  \end{prooftree}
  \\
  \tag{Domain widens}
  \begin{prooftree}
    \hypo{a: T[domain=m]}
    \hypo{b: T[domain=n]}
    \hypo{T \in \{Int, Float\}}
    \infer3{(a + b): T[domain=map(sum, m \times n)]}
  \end{prooftree}
\end{gather}
\subparagraph{Subtract}
\begin{gather}
  \tag{Literal}
  \begin{prooftree}
    \hypo{a: T[literal=x]}
    \hypo{b: T[literal=y]}
    \hypo{T \in \{Int, Float\}}
    \infer3{(a - b): T[literal=(x - y)]}
  \end{prooftree}
  \\
  \tag{Literal and min}
  \begin{prooftree}
    \hypo{a: T[literal=x]}
    \hypo{b: T[min=y]}
    \hypo{T \in \{Int, Float\}}
    \infer3{(a - b): T[max=(x - y)]}
  \end{prooftree}
  \\
  \tag{Literal and max}
  \begin{prooftree}
    \hypo{a: T[literal=x]}
    \hypo{b: T[max=y]}
    \hypo{T \in \{Int, Float\}}
    \infer3{(a - b): T[min=(x - y)]}
  \end{prooftree}
  \\
  \tag{Min and literal}
  \begin{prooftree}
    \hypo{a: T[min=x]}
    \hypo{b: T[literal=y]}
    \hypo{T \in \{Int, Float\}}
    \infer3{(a - b): T[min=(x-y)]}
  \end{prooftree}
  \\
  \tag{Max and literal}
  \begin{prooftree}
    \hypo{a: T[max=x]}
    \hypo{b: T[literal=y]}
    \hypo{T \in \{Int, Float\}}
    \infer3{(a - b): T[max=(x-y)]}
  \end{prooftree}
  \\
  \tag{Domain widens}
  \begin{prooftree}
    \hypo{a: T[domain=m]}
    \hypo{b: T[domain=n]}
    \hypo{T \in \{Int, Float\}}
    \infer3{(a - b): T[domain=map(sum, m \times map(-, n))]}
  \end{prooftree}
\end{gather}
\subparagraph{Divide}
\begin{gather}
  \tag{Literal}
  \begin{prooftree}
    \hypo{a: Int[literal=x]}
    \hypo{b: Int[literal=y]}
    \infer2{(a / b): T[literal=(x / y)]}
  \end{prooftree}
  \\
  \tag{Literal zero}
  \begin{prooftree}
    \hypo{a: Int[literal=0]}
    \hypo{b: Int[0 \not\in domain \lor max < 0 \lor min > 0]}
    \hypo{x \geq 0}
    \infer3{(a / b): T[literal=0]}
  \end{prooftree}
  \\
  \tag{Min and literal}
  \begin{prooftree}
    \hypo{a: Int[min=x]}
    \hypo{b: Int[literal=y]}
    \hypo{x \geq 0}
    \infer3{(a / b): T[min=(x / y)]}
  \end{prooftree}
  \\
  \tag{Max and literal}
  \begin{prooftree}
    \hypo{a: Int[max=x]}
    \hypo{b: Int[literal=y]}
    \infer2{(a / b): T[max=(x / y)]}
  \end{prooftree}
  \\
  \tag{Literal and min}
  \begin{prooftree}
    \hypo{a: Int[literal=x]}
    \hypo{b: Int[min=y]}
    \hypo{x > 0}
    \hypo{y > 0}
    \infer4{(a / b): T[max=(x / y)]}
  \end{prooftree}
  \\
  \tag{Literal and max}
  \begin{prooftree}
    \hypo{a: Int[literal=x]}
    \hypo{b: Int[max=y]}
    \hypo{x > 0}
    \hypo{y < 0}
    \infer4{(a / b): T[min=(x / y)]}
  \end{prooftree}
  \\
  \tag{Negative literal and min}
  \begin{prooftree}
    \hypo{a: Int[literal=x]}
    \hypo{b: Int[min=y]}
    \hypo{x < 0}
    \hypo{y > 0}
    \infer4{(a / b): T[min=(x / y)]}
  \end{prooftree}
  \\
  \tag{Negative literal and max}
  \begin{prooftree}
    \hypo{a: Int[literal=x]}
    \hypo{b: Int[max=y]}
    \hypo{x < 0}
    \hypo{y < 0}
    \infer4{(a / b): T[max=(x / y)]}
  \end{prooftree}
\end{gather}
\subparagraph{Multiply} And commutative versions.
\begin{gather}
  \tag{Literal}
  \begin{prooftree}
    \hypo{a: T[literal=x]}
    \hypo{b: T[literal=y]}
    \hypo{T \in \{Int, Float\}}
    \infer3{(a \times b): T[literal=(x \times y)]}
  \end{prooftree}
  \\
  \tag{Literal zero}
  \begin{prooftree}
    \hypo{a: T[literal=0]}
    \hypo{b: T}
    \hypo{T \in \{Int, Float\}}
    \infer3{(a \times b): T[literal=0]}
  \end{prooftree}
  \\
  \tag{Positive literal and min}
  \begin{prooftree}
    \hypo{a: T[literal=x]}
    \hypo{b: T[min=y]}
    \hypo{T \in \{Int, Float\}}
    \hypo{x > 0}
    \infer4{(a \times b): T[min=(x \times y)]}
  \end{prooftree}
  \\
  \tag{Positive literal and max}
  \begin{prooftree}
    \hypo{a: T[literal=x]}
    \hypo{b: T[max=y]}
    \hypo{T \in \{Int, Float\}}
    \hypo{x > 0}
    \infer4{(a \times b): T[max=(x \times y)]}
  \end{prooftree}
  \\
  \tag{Negative literal and min}
  \begin{prooftree}
    \hypo{a: T[literal=x]}
    \hypo{b: T[min=y]}
    \hypo{T \in \{Int, Float\}}
    \hypo{x < 0}
    \infer4{(a \times b): T[max=(x \times y)]}
  \end{prooftree}
  \\
  \tag{Negative literal and max}
  \begin{prooftree}
    \hypo{a: T[literal=x]}
    \hypo{b: T[max=y]}
    \hypo{T \in \{Int, Float\}}
    \hypo{x < 0}
    \infer4{(a \times b): T[max=(x \times y)]}
  \end{prooftree}
  \\
  \tag{Domain widens}
  \begin{prooftree}
    \hypo{a: T[domain=m]}
    \hypo{b: T[domain=n]}
    \hypo{T \in \{Int, Float\}}
    \infer3{(a \times b): T[domain=map(product, m \times n)]}
  \end{prooftree}
\end{gather}
\subparagraph{Power}
\begin{gather}
  \tag{Literal}
  \begin{prooftree}
    \hypo{a: T[literal=x]}
    \hypo{b: T[literal=y]}
    \hypo{T \in \{Int, Float\}}
    \infer3{(a^b): T[literal=(x^y)]}
  \end{prooftree}
  \\
  \tag{Literal zero}
  \begin{prooftree}
    \hypo{a: T[literal=0]}
    \hypo{b: T[literal != 0]}
    \hypo{T \in \{Int, Float\}}
    \infer3{(a^b): T[literal=0]}
  \end{prooftree}
  \\
  \tag{Literal zero}
  \begin{prooftree}
    \hypo{a: T}
    \hypo{b: T[literal=0]}
    \hypo{T \in \{Int, Float\}}
    \infer3{(a^b): T[literal=1]}
  \end{prooftree}
  \\
  \tag{Literal and min}
  \begin{prooftree}
    \hypo{a: T[literal=x]}
    \hypo{b: T[min=y]}
    \hypo{T \in \{Int, Float\}}
    \infer3{(a^b): T[min=(x^y)]}
  \end{prooftree}
  \\
  \tag{Literal and max}
  \begin{prooftree}
    \hypo{a: T[literal=x]}
    \hypo{b: T[max=y]}
    \hypo{T \in \{Int, Float\}}
    \infer3{(a^b): T[max=(x^y)]}
  \end{prooftree}
  \\
  \tag{Domain widens}
  \begin{prooftree}
    \hypo{a: T[domain=m]}
    \hypo{b: T[domain=n]}
    \hypo{T \in \{Int, Float\}}
    \infer3{(a^b): T[domain=map(\lambda (x,y):x^y, m \times n)]}
  \end{prooftree}
\end{gather}
\subparagraph{Upto}
\begin{gather}
  \tag{Upto literal}
  \begin{prooftree}
    \hypo{a: Int[literal=x]}
    \hypo{b: Int[literal=y]}
    \infer2{(a \upto b): set[min=x, max=y, literal=\{x, x+1, ..., y\}, total]}
  \end{prooftree} % If only one of these are literals, we should still try to get either a max or min value
  \\
  \tag{Upto literal}
  \begin{prooftree}
    \hypo{a: Int[min=x]}
    \hypo{b: Int}
    \infer2{(a \upto b): set[min=x]}
  \end{prooftree}
  \\
  \tag{Upto literal}
  \begin{prooftree}
    \hypo{a: Int}
    \hypo{b: Int[max=x]}
    \infer2{(a \upto b): set[max=x]}
  \end{prooftree}
\end{gather}
% \paragraph{String}
\paragraph{Tuple}
\subparagraph{Enumeration}
\begin{gather}
  \tag{Literal}
  \begin{prooftree}
    \hypo{a_1: T_1, a_2: T_2, ..., a_n: T_n}
    \infer1{(a_1, a_2, ..., a_n): tuple[T_1, T_2, ..., T_n, literal=(a_1, a_2, ..., a_n), size=n]}
  \end{prooftree}
\end{gather}
\subparagraph{Maplet}
\begin{gather}
  \tag{Literal}
  \begin{prooftree}
    \hypo{a_1: T_1, a_2: T_2}
    \infer1{(a_1 \mapsto a_2): tuple[T_1, T_2, literal=(a_1, a_2), size=2]}
  \end{prooftree}
\end{gather}
% \paragraph{Record}
% \subparagraph{Access}
\paragraph{Set}
Note: the domain trait on sets is not the same as that of primitive types. A primitive value selects one out of the domain, whereas collection types choose from a subset of the domain. For example, let $x: int[domain=\{1,2,3\}]$. Then $x \in \{1,2,3\}$. However, if $y: set[int, domain=\{1,2,3\}]$, then $y \subset \{1,2,3\}$ (that is, $y \in \mathbb{P}\{1,2,3\}$).
\subparagraph{Enumeration}
\begin{gather}
  \tag{Literal}
  \begin{prooftree}
    \hypo{a_1, a_2, ..., a_n: T}
    \infer1{\{a_1, a_2, ..., a_n\}: set[T, literal=\{a_1, a_2, ..., a_n\}, size=n]}
  \end{prooftree}
\end{gather}
\subparagraph{Clear}
\begin{gather}
  \tag{Literal}
  \begin{prooftree}
    \hypo{s: set[T]}
    \infer1{clear(s): set[T, literal=\{\}, empty]}
  \end{prooftree}
\end{gather}
\subparagraph{Is empty}
\begin{gather}
  \tag{Literal}
  \begin{prooftree}
    \hypo{s: set[Empty]}
    \infer1{(s = \{\}): Bool[literal=true]}
  \end{prooftree}
\end{gather}
\subparagraph{Add}
\begin{gather}
  \tag{Size may increase}
  \begin{prooftree}
    \hypo{s: set[T, Size=x]}
    \hypo{a: T}
    \infer2{(s \cup \{a\}): set[T, Size=x+1]} % Can check literals to see if the size will actually increase. For now, best overestimation guess
  \end{prooftree}
  \\
  \tag{Total size stays the same}
  \begin{prooftree}
    \hypo{s: set[T, Size=x, Total]}
    \hypo{a: T}
    \infer2{(s \cup \{a\}): set[T, Size=x, Total]}
  \end{prooftree}
  \\
  \tag{Literal}
  \begin{prooftree}
    \hypo{s: set[T, literal=l_1]}
    \hypo{a: T[literal=l_2]}
    \infer2{(s \cup \{a\}): set[T, literal=(l_1 \cup \{l_2\})]}
  \end{prooftree}
  \\
  \tag{Domain with literal}
  \begin{prooftree}
    \hypo{s: set[T, domain=d]}
    \hypo{a: T[literal=l]}
    \infer2{(s \cup \{a\}): set[T, domain=(d \cup \{l\})]}
  \end{prooftree}
  \\
  \tag{Domain with domain}
  \begin{prooftree}
    \hypo{s: set[T, domain=d]}
    \hypo{a: T[domain=l]}
    \infer2{(s \cup \{a\}): set[T, literal=(d \cup l)]}
  \end{prooftree}
  \\
  \tag{Empty}
  \begin{prooftree}
    \hypo{s: set[T, empty]}
    \hypo{a: T}
    \infer2{(s \cup \{a\}): set[T]}
  \end{prooftree}
\end{gather}
\subparagraph{Remove}
\begin{gather}
  \tag{Size does not guarantee a decrease}
  \begin{prooftree}
    \hypo{s: set[T, size=n]}
    \hypo{a: T}
    \infer2{(s \setminus \{a\}): set[T, size=n]}
  \end{prooftree}
  \\
  \tag{Removes Total guarantee}
  \begin{prooftree}
    \hypo{s: set[T, size=x, total]}
    \hypo{a: T}
    \infer2{(s \setminus \{a\}): set[T, size=x-1]}
  \end{prooftree}
  \\
  \tag{Literal}
  \begin{prooftree}
    \hypo{s: set[T, literal=l_1]}
    \hypo{a: T[literal=l_2]}
    \infer2{(s \setminus \{a\}): set[T, literal=(l_1 \setminus \{l_2\})]}
  \end{prooftree}
  \\
  \tag{Domain with literal}
  \begin{prooftree}
    \hypo{s: set[T, domain=d]}
    \hypo{a: T[literal=l]}
    \infer2{(s \setminus \{a\}): set[T, domain=(d \setminus \{l\})]}
  \end{prooftree}
  \\
  \tag{Empty}
  \begin{prooftree}
    \hypo{s: set[T, empty]}
    \hypo{a: T}
    \infer2{(s \setminus \{a\}): set[T, empty]}
  \end{prooftree}
\end{gather}
\subparagraph{In} Not in is similar.
\begin{gather}
  \tag{Literal}
  \begin{prooftree}
    \hypo{s: set[T, literal=l_1]}
    \hypo{a: T[literal=l_2]}
    \infer2{(a \in s): Bool[literal=(l_2 \in l_1)]}
  \end{prooftree}
  \\
  \tag{Total domain with literal}
  \begin{prooftree}
    \hypo{s: set[T, domain=d, total]}
    \hypo{a: T[literal=l]}
    \infer2{(a \in s): Bool[literal=(l \in d)]}
  \end{prooftree}
  \\
  \tag{Total domain with domain}
  \begin{prooftree}
    \hypo{s: set[T, domain=d_1, total]}
    \hypo{a: T[domain=d_2]}
    \infer2{(a \in s): Bool[literal=(d_2 \subseteq d_1)]}
  \end{prooftree}
  \\
  \tag{Empty}
  \begin{prooftree}
    \hypo{s: set[T, empty]}
    \hypo{a: T}
    \infer2{(a \in s): Bool[literal=false]}
  \end{prooftree}
\end{gather}
\subparagraph{Cardinality}
\begin{gather}
  \tag{Size overestimates}
  \begin{prooftree}
    \hypo{s: set[T, size=n]}
    \infer1{card(s): Int[max=n]}
  \end{prooftree}
  \\
  \tag{Literal}
  \begin{prooftree}
    \hypo{s: set[T, literal=l]}
    \infer1{card(s): Int[literal=card(l)]}
  \end{prooftree}
  \\
  \tag{Domain}
  \begin{prooftree}
    \hypo{s: set[T, domain=d]}
    \infer1{card(s): Int[max=card(d)]}
  \end{prooftree}
  \\
  \tag{Empty}
  \begin{prooftree}
    \hypo{s: set[T, empty]}
    \infer1{card(s): Int[literal=0]}
  \end{prooftree}
\end{gather}
\subparagraph{Powerset}
\begin{gather}
  \tag{Total size}
  \begin{prooftree}
    \hypo{s: set[T, size=n, total]}
    \infer1{powerset(s): set[set[T], size=2^n, total]}
  \end{prooftree}
  \\
  \tag{Literal}
  \begin{prooftree}
    \hypo{s: set[T, literal=l]}
    \infer1{powerset(s): set[set[T], literal=\mathbb{P}(l)]}
  \end{prooftree}
  \\
  \tag{Domain}
  \begin{prooftree}
    \hypo{s: set[T, domain=d]}
    \infer1{powerset(s): set[set[T], domain=\mathbb{P}(d)]}
  \end{prooftree}
\end{gather}
\subparagraph{Choice}
\begin{gather}
  \tag{Empty choice fails (statically)}
  \begin{prooftree}
    \hypo{s: set[T, empty]}
    \infer1{choice(s): T[undefined]}
  \end{prooftree}
  \\
  \tag{Single literal}
  \begin{prooftree}
    \hypo{s: set[T, literal=\{l\}]}
    \infer1{choice(s): T[literal=l]}
  \end{prooftree}
  \\
  \tag{Domain}
  \begin{prooftree}
    \hypo{s: set[T, domain=d]}
    \infer1{choice(s): T[domain=l]}
  \end{prooftree}
\end{gather}
\subparagraph{Min}
\begin{gather}
  \tag{Min}
  \begin{prooftree}
    \hypo{s: set[T, min=m]}
    \infer1{min(s): T[literal=m]}
  \end{prooftree}
  \\
  \tag{Single literal}
  \begin{prooftree}
    \hypo{s: set[T, literal=\{l\}]}
    \infer1{min(s): T[literal=l]}
  \end{prooftree}
  \\
  \tag{Literal}
  \begin{prooftree}
    \hypo{s: set[T, literal=l]}
    \infer1{min(s): T[domain=l]}
  \end{prooftree}
  \\
  \tag{Domain}
  \begin{prooftree}
    \hypo{s: set[T, domain=d]}
    \infer1{min(s): T[domain=d]}
  \end{prooftree}
\end{gather}
\subparagraph{Max}
\begin{gather}
  \tag{Max}
  \begin{prooftree}
    \hypo{s: set[T, max=m]}
    \infer1{max(s): T[literal=m]}
  \end{prooftree}
  \\
  \tag{Single literal}
  \begin{prooftree}
    \hypo{s: set[T, literal=\{l\}]}
    \infer1{max(s): T[literal=l]}
  \end{prooftree}
  \\
  \tag{Literal}
  \begin{prooftree}
    \hypo{s: set[T, literal=l]}
    \infer1{max(s): T[domain=l]}
  \end{prooftree}
  \\
  \tag{Domain}
  \begin{prooftree}
    \hypo{s: set[T, domain=d]}
    \infer1{max(s): T[domain=d]}
  \end{prooftree}
\end{gather}
\subparagraph{Sum}
\begin{gather}
  \tag{Sum bounds}
  \begin{prooftree}
    \hypo{s: set[T, max=y, size=n]}
    \infer1{sum(s): T[max=y\times n]}
  \end{prooftree}
  \\
  \tag{Literal}
  \begin{prooftree}
    \hypo{s: set[T, literal=l]}
    \infer1{sum(s): T[literal=sum(l)]}
  \end{prooftree}
\end{gather}
\subparagraph{Product}
\begin{gather}
  \tag{Product bounds}
  \begin{prooftree}
    \hypo{s: set[T, max=y, size=n]}
    \infer1{product(s): T[max=y^n]}
  \end{prooftree}
  \\
  \tag{Literal}
  \begin{prooftree}
    \hypo{s: set[T, literal=l]}
    \infer1{product(s): T[literal=product(l)]}
  \end{prooftree}
\end{gather}
\subparagraph{Union}
\begin{gather}
  \tag{Literal}
  \begin{prooftree}
    \hypo{s: set[T, literal=l_1]}
    \hypo{t: set[T, literal=l_2]}
    \infer2{(s \cup t): set[T, literal=(l_1 \cup l_2)]}
  \end{prooftree}
  \\
  \tag{Domain}
  \begin{prooftree}
    \hypo{s: set[T, domain=d_1]}
    \hypo{t: set[T, domain=d_2]}
    \infer2{(s \cup t): set[T, domain=(d_1 \cup d_2)]}
  \end{prooftree}
  \\
  \tag{Domain and Literal}
  \begin{prooftree}
    \hypo{s: set[T, domain=d]}
    \hypo{t: set[T, literal=l]}
    \infer2{(s \cup t): set[T, domain=(d \cup l)]}
  \end{prooftree}
  \\
  \tag{Lowest Min}
  \begin{prooftree}
    \hypo{s: set[T, min=m_1]}
    \hypo{t: set[T, min=m_2]}
    \hypo{m_1 < m_2}
    \infer3{(s \cup t): set[T, min=m_1]}
  \end{prooftree}
  \\
  \tag{Highest Max}
  \begin{prooftree}
    \hypo{s: set[T, max=m_1]}
    \hypo{t: set[T, max=m_2]}
    \hypo{m_1 < m_2}
    \infer3{(s \cup t): set[T, max=m_2]}
  \end{prooftree}
  \\
  \tag{Union with Empty keeps Traits}
  \begin{prooftree}
    \hypo{s: set[T, empty]}
    \hypo{t: set[T, traits]}
    \infer2{(s \cup t): set[T, traits]}
  \end{prooftree}
  \\
  \tag{Size overestimates}
  \begin{prooftree}
    \hypo{s: set[T, size=x]}
    \hypo{t: set[T, size=y]}
    \infer2{(s \cup t): set[T, size=(x+y)]}
  \end{prooftree}
  \\
  \tag{Total Domain}
  \begin{prooftree}
    \hypo{s: set[T, domain=d_1]}
    \hypo{t: set[T, domain=d_2, total]}
    \hypo{d_1 \subseteq d_2}
    \infer3{(s \cup t): set[T, domain=d_2, total]}
  \end{prooftree}
\end{gather}
\subparagraph{Intersection}
\begin{gather}
  \tag{Literal}
  \begin{prooftree}
    \hypo{s: set[T, literal=l_1]}
    \hypo{t: set[T, literal=l_2]}
    \infer2{(s \cap t): set[T, literal=(l_1 \cap l_2)]}
  \end{prooftree}
  \\
  \tag{Domain}
  \begin{prooftree}
    \hypo{s: set[T, domain=d_1]}
    \hypo{t: set[T, domain=d_2]}
    \infer2{(s \cap t): set[T, domain=(d_1 \cap d_2)]}
  \end{prooftree}
  \\
  \tag{Domain and Literal}
  \begin{prooftree}
    \hypo{s: set[T, domain=d]}
    \hypo{t: set[T, literal=l]}
    \infer2{(s \cap t): set[T, domain=(d \cap l)]}
  \end{prooftree}
  \\
  \tag{Highest Min}
  \begin{prooftree}
    \hypo{s: set[T, min=m_1]}
    \hypo{t: set[T, min=m_2]}
    \hypo{m_1 < m_2}
    \infer3{(s \cap t): set[T, min=m_2]}
  \end{prooftree}
  \\
  \tag{Lowest Max}
  \begin{prooftree}
    \hypo{s: set[T, max=m_1]}
    \hypo{t: set[T, max=m_2]}
    \hypo{m_1 < m_2}
    \infer3{(s \cap t): set[T, max=m_1]}
  \end{prooftree}
  \\
  \tag{Intersection with Empty is Empty}
  \begin{prooftree}
    \hypo{s: set[T, empty]}
    \hypo{t: set[T]}
    \infer2{(s \cap t): set[T, empty]}
  \end{prooftree}
  \\
  \tag{Size does not increase}
  \begin{prooftree}
    \hypo{s: set[T, size=x]}
    \hypo{t: set[T, size=y]}
    \infer2{(s \cap t): set[T, size=max(x,y)]}
  \end{prooftree}
  \\
  \tag{Total Domain}
  \begin{prooftree}
    \hypo{s: set[T, domain=d_1]}
    \hypo{t: set[T, domain=d_2, total]}
    \hypo{d_1 \subseteq d_2}
    \infer3{(s \cup t): set[T, domain=d_1]}
  \end{prooftree}
\end{gather}
\subparagraph{Difference}
\begin{gather}
  \tag{Literal}
  \begin{prooftree}
    \hypo{s: set[T, literal=l_1]}
    \hypo{t: set[T, literal=l_2]}
    \infer2{(s \setminus t): set[T, literal=(l_1 \setminus l_2)]}
  \end{prooftree}
  \\
  \tag{Literal}
  \begin{prooftree}
    \hypo{s: set[T, literal=l]}
    \hypo{t: set[T]}
    \infer2{(s \setminus t): set[T, domain=l]}
  \end{prooftree}
  \\
  \tag{Domain}
  \begin{prooftree}
    \hypo{s: set[T, domain=d]}
    \hypo{t: set[T]}
    \infer2{(s \setminus t): set[T, domain=d]}
  \end{prooftree}
  \\
  \tag{Domain and Literal}
  \begin{prooftree}
    \hypo{s: set[T, domain=d]}
    \hypo{t: set[T, literal=l]}
    \infer2{(s \setminus t): set[T, domain=(d \setminus l)]}
  \end{prooftree}
  \\
  \tag{Difference with Empty keeps traits}
  \begin{prooftree}
    \hypo{s: set[T, traits]}
    \hypo{t: set[T, empty]}
    \infer2{(s \setminus t): set[T, traits]}
  \end{prooftree}
  \\
  \tag{Empty difference}
  \begin{prooftree}
    \hypo{s: set[T, empty]}
    \hypo{t: set[T]}
    \infer2{(s \setminus t): set[T, empty]}
  \end{prooftree}
  \\
  \tag{Size does not increase}
  \begin{prooftree}
    \hypo{s: set[T, size=x]}
    \hypo{t: set[T]}
    \infer2{(s \setminus t): set[T, size=x]}
  \end{prooftree}
\end{gather}
% \subparagraph{Symmetric difference}
\subparagraph{Cartesian product}
\begin{gather}
  \tag{Literal}
  \begin{prooftree}
    \hypo{s: set[T, literal=l_1]}
    \hypo{t: set[V, literal=l_2]}
    \infer2{(s \times t): relation[T, V, literal=l_1 \times l_2]}
  \end{prooftree}
  \\
  \tag{Domain}
  \begin{prooftree}
    \hypo{s: set[T, domain=d_1]}
    \hypo{t: set[V, domain=d_2]}
    \infer2{(s \times t): relation[T, V, domain=d_1 \times d_2]}
  \end{prooftree}
  \\
  \tag{Domain and Literal}
  \begin{prooftree}
    \hypo{s: set[T, domain=d]}
    \hypo{t: set[V, literal=l]}
    \infer2{(s \times t): relation[T, V, domain=d \times l]}
  \end{prooftree}
  \\
  \tag{Literal and Domain}
  \begin{prooftree}
    \hypo{s: set[T, literal=l]}
    \hypo{t: set[V, domain=d]}
    \infer2{(s \times t): relation[T, V, literal=l \times d]}
  \end{prooftree}
  \\
  \tag{Total}
  \begin{prooftree}
    \hypo{s: set[T, domain=d_1, total]}
    \hypo{t: set[V, domain=d_2, total]}
    \infer2{(s \times t): relation[T, V, domain=d_1 \times d_2, total]}
  \end{prooftree}
  \\
  \tag{Size}
  \begin{prooftree}
    \hypo{s: set[T, size=x]}
    \hypo{t: set[V, size=y]}
    \infer2{(s \times t): relation[T, V, size=x \times y]}
  \end{prooftree}
  \\
  \tag{Empty}
  \begin{prooftree}
    \hypo{s: set[T, empty]}
    \hypo{t: set[V]}
    \infer2{(s \times t): relation[T, V, empty]}
  \end{prooftree}
  \\
  \tag{Empty}
  \begin{prooftree}
    \hypo{s: set[T]}
    \hypo{t: set[V, empty]}
    \infer2{(s \times t): relation[T, V, empty]}
  \end{prooftree}
\end{gather}
% \subparagraph{Disjoint} % What operator is this?
\subparagraph{Subset}
\begin{gather}
  \tag{Subset with min-max}
  \begin{prooftree}
    \hypo{s: set[T, Min=w, Max=x]}
    \hypo{t: set[T, Min=y, Max=z]}
    \hypo{w > y}
    \hypo{x < z}
    % Then we should only consider values within w and x range? Assuming they are ordered...
    \infer4{(s \subset t: set[T, Min=w, Max=x]): Bool}
  \end{prooftree}
\end{gather}
\subparagraph{Subset equals}
\subparagraph{Superset}
\subparagraph{Superset equals}
\subparagraph{Not subset}
\subparagraph{Not subset equals}
\subparagraph{Not superset}
\subparagraph{Not superset equals}
\paragraph{Relation}
\subparagraph{Inverse}
\begin{gather}
  % NOTE: These are in addition to the relation-specific traits in the above section!
  \tag{Inverse swaps Domain}
  \begin{prooftree}
    \hypo{r: relation[T[Domain=x], V[Domain=y]]}
    \infer1{r^{-1}: relation[V[Domain=y], T[Domain=x]]}
  \end{prooftree}
\end{gather}
\subparagraph{Composition}
\begin{gather}
   % NOTE: These are in addition to the relation-specific traits in the above section!
  \tag{Empty Composition}
  \begin{prooftree}
    \hypo{s: relation[T, V, Empty]}
    \hypo{t: relation[V, U]} % this one may be empty too
    \infer2{s \circ t: relation[T, U, Empty]}
  \end{prooftree}
\end{gather}
\subparagraph{Function call}
\subparagraph{Image}
\begin{gather}
  % NOTE: These are in addition to the relation-specific traits in the above section!
  \tag{Image with Empty Arg}
  \begin{prooftree}
    \hypo{r: relation[T, V]}
    \hypo{s: set[T, Empty]}
    \infer2{r[s]: set[V, Empty]}
  \end{prooftree}
  \\
  % NOTE: These are in addition to the relation-specific traits in the above section!
  \tag{Image with Empty Relation}
  \begin{prooftree}
    \hypo{r: relation[T, V, Empty]}
    \hypo{s: set[T]}
    \infer2{r[s]: set[V, Empty]}
  \end{prooftree}
  % TODO more can be determined to be empty by using max/min/domain
  \\
  % NOTE: These are in addition to the relation-specific traits in the above section!
  \tag{Image with ManyToOne}
  \begin{prooftree}
    \hypo{r: relation[T, V, TotalOnDomain]}
    \hypo{s: set[T, Domain=domain(r), Size > 0]}
    \infer2{r[s]: set[V, Size > 0]}
  \end{prooftree}
  \\
  % NOTE: These are in addition to the relation-specific traits in the above section!
  \tag{Image with ManyToOne}
  \begin{prooftree}
    \hypo{r: relation[T, V, ManyToOne]}
    \hypo{s: set[T, Size = n]}
    \infer2{r[s]: set[V, Size = n]}
  \end{prooftree}
\end{gather}
\subparagraph{Overriding}
\subparagraph{Domain}
\begin{gather}
  \tag{Total domain}
  \begin{prooftree}
    \hypo{r: relation[T, V, TotalOnDomain]}
    \infer1{domain(r): set[T, Total]}
  \end{prooftree}
\end{gather}
\subparagraph{Range}
\begin{gather}
  \tag{Total range}
  \begin{prooftree}
    \hypo{r: relation[T, V, TotalOnRange]}
    \infer1{range(r): set[V, Total]}
  \end{prooftree}
\end{gather}
\subparagraph{Domain restriction}
\subparagraph{Domain subtraction}
\subparagraph{Range restriction}
\subparagraph{Range subtraction}
\subparagraph{Bag Image}
\subparagraph{Comprehension}
\paragraph{Bag}
\subparagraph{Bag union}
\subparagraph{Bag intersection}
\subparagraph{Bag add}
\subparagraph{Bag difference}
\subparagraph{Comprehension}
\subparagraph{Size}
\paragraph{Sequence}
\subparagraph{Concat}
\begin{gather}
  \tag{Concat with Size}
  \begin{prooftree}
    \hypo{s: sequence[T, Size=x]}
    \hypo{t: sequence[T, Size=y]}
    \infer2{s \concat t: sequence[T, Size=x+y]}
  \end{prooftree}
\end{gather}
\subparagraph{Comprehension}
\paragraph{Procedure}
\subparagraph{Call}
    % Procedure call? need to propagate traits from return value type - should be done at procedure definition time?
\paragraph{Quantification Body}
\subparagraph{Min}
\subparagraph{Max}
\subparagraph{Forall}
\subparagraph{Exists}
\subparagraph{Union all}
\subparagraph{Intersection all}
\subparagraph{Comprehension}
\subparagraph{Sum}
\subparagraph{Product}
\subparagraph{Map}
\subparagraph{Map min}
\subparagraph{Map max}
\subparagraph{Iter}
\subparagraph{Fold}

% \begin{gather}
%     \tag{Generics with set ops}
%     \begin{prooftree}
%         \hypo{s: set[T[GenericBound=x]]}
%         \hypo{t: set[T[GenericBound=y]]}
%         \infer2{(s \oplus t): set[T[GenericBound=x \cap y]]}
%     \end{prooftree}
%     % Intersection and difference, much of the same but take minimum, order matters for difference (eg difference from empty set), precompute literals, set operations takes the minimum of bound types for generics
% \end{gather}

% TODO add a section for value-trait interactions (eg. errors/undefinedness, literal assignment, etc.)

\subsection{Definedness}
% Operators that can be undefined:
% - division by 0
% - function on an empty set, simplifies to choice of image
% - all non-equality operations with undefined as an input is undefined
% - complex numbers?
% - 0^0?
% the event b guide specifies short circuited ands for definedness, but they probably dont swap ands in optimization... we should try two different ways: either all and clauses must be defined or we attempt to recover from undefined clauses when seen
% Functions that can be undefined:
% - choice on empty set (undefined with a given type)

\subsection{Language Examples} % Move to motivation?
\section{Features and Behaviours}



\section{(Abstract) Data Types}

% Table to summarize types, traits, and operations

\subsection{Traits}
% Description, Purpose, and Motivation: What is this construct for? Why do we need it? What is it based off of (find source)?
% Semantics: What do traits represent? What does it mean when a type has a trait?
% For each trait:
% - Description, applications
% - Operations affected/unlocked by this trait
% - Assumptions, Invariants/Laws
% Maybe a table to summarize each trait?

\subsection{Boolean}
\subsection{Integer}
\subsection{Float}
\subsection{String}
\subsection{Set}
\subsection{Relation}
\subsection{Sequence}
\subsection{Bag}
\subsection{Tuple}
\subsection{Record}
\subsection{Procedure}
\subsection{Generic}
% Description, Purpose, and Motivation: What is this type for? Why do we need it? What is it based off of (find source)?
% Semantics: What do values represent? What does it mean to have an object of this type?
% Making this type: Example of its construction/constructors
% Examples of its use
% Relationships to other types (super/subtypes)
% Allowed traits (and how they affect this type)
% Restricted traits (and why we restrict them; what are the alternatives?)
% Formal definition
% Assumptions, Invariants/Laws
% Any other notes (edge cases, differences from normal execution)

\subsubsection{Operations}
% For each operation (include basic operations like equality, etc.):
% - Description: What is it used for? What is its purpose? Why do we need it? Where does it come from (find source)?
% - Example usage
% - Traits that affect this operator (and how they affect the operator)
% - Formalized type definition (with and without traits)
% - Any other notes/Detailed knowledge on this operator (edge cases, differences from normal execution)
% - Assumptions, Invariants/Laws















\section{User-Facing Specification}
% Non-determinism, short circuiting

\subsection{(Abstract) Data Types}
For the most part, Simile's semantics resemble those of set theory and specification languages (specifically, Event-B). Aside from a few primitive types that abstract over bare set theory and programming language specific structures, like integer arithmetic and procedure definitions, all objects inherit properties from sets. For example, a sequence can be modelled as a set of pairs from natural numbers to the sequence element type.

In the forthcoming sections, we explicitly state the design of these set-theory-inspired types, useful properties, and underlying implementations.

\subsubsection{Primitives}
With consideration to execution efficiency and the extensive amount of existing optimization in modern languages, Simile makes use of the following basic types:
\begin{description}
    \item[Booleans] As the cornerstone of logical reasoning, booleans deserve an explicit type with reserved identifiers: $\mathbb{B} = \{True, False\}$. Further, Simile provides a healthy supply of common operations beyond $\land$ and $\lor$: $\implies, \equiv, \impliedby, \forall, \exists,$ etc.
    \item[Integers] While integers are a very mature type in the world of computing, trade-offs between arbitrary-precision and fixed-point arithmetic may require testing and careful review.
    \\
    If Simile chooses to represent numbers through conventional fixed point arithmetic, struggles with overflow and extremely large numbers are compounded by the innate non-determinism within Simile. For example, calculating the sum of a set of very large numbers can be done in arbitrary order, according to the semantics of mathematics (and Simile). But in fixed-point arithmetic, the resulting value can very well differ based on the order of operations. In other words, commutative math operations are no longer commutative under fixed precision. At some point in the computation pipeline, Simile will need to decide the order in which to evaluate the set, which may cause hard-to-spot bugs and unexpected runtime failures.
    \\
    Conversely, choosing arbitrary precision integers may bring Simile operators closer to mathematics at the cost of speed. The core features of Simile, while targeting formal methods users, attempt to efficiently execute set and relation operations. Adding checks for overflow on every mathematical operation may prove detrimental to performance.
    \\
    The implementation decision essentially is one of context vs. safety. To benchmark Simile with these types, we gather a list of examples and judge whether the performance hit is worth the reliability benefits. Specific questions to examine: How often does overflow occur in these examples? Do numbers approach the upper limits of fixed precision often enough to warrant accommodating infrastructure? What is the performance of fixed vs. arbitrary precision calculations.
    \item[Floating Point Numbers] In a similar vein to integers, floats are an imperfect translation of real numbers into efficiently calculable representations. Since floats (or even reals) are not used as often in specifications as integers, we leave the implementation as the conventional 64-bit representation. However, we note that non-determinism on these operations will prove especially tricky, since the result of commutative complements may differ without explicit over/underflow.
    \item[Strings] Strings are syntactic sugar over Sequences of characters, represented internally as a sequence of (fixed precision) integers. Eventually, Simile aims to support built-in string operations as seen in many common programming languages.
\end{description}

\subsubsection{Sets}
Foundational Simile type: an unordered, unique collection of objects.

Sets may be used in two ways:
\begin{enumerate}
    \item Enumerations by way of $\{x,y,z\}$ (read: The set of elements $x,y$, and $z$)
    \item Comprehensions of the form $\bSet{x}{x \in S}{f(x)}$ (read: The set of function $f$ applied to all elements in $S$). A small effort is made to accommodate a shorthand form, wherein $\bSet{x}{x \in S}{x}$ may be written as $\bSetT{x}{x \in S}$, but this syntax easily confuses bound/unbound variables and is not recommended for complex comprehension expressions.
\end{enumerate}

Sets, as mathematical objects with only two major restrictions, lend themselves well to a variety of computational implementations. Arrays, ordered arrays, hash maps, bit sets, roaring bit sets, path maps, and tries can all be used to satisfy the uniqueness property and hide an internally ordered representation from users.

% Arrays, heaps, hashmaps, bitsets, roaring bitsets, pathmaps, tries?

No matter the underlying implementation, the optimizer may make use of syntactically-analyzed set sizes to direct the choice of lowered iterator. Further, recording the cardinality of a set at runtime may prove beneficial to the performance of common cardinality-based specification expressions at a very small memory cost. Other properties based on element type, like max/min for numeric types, limits for numeric ranges, maximum sizes, may be useful.


\subsubsection{Relations}
Relations are sets that make use of paired/maplet elements (i.e., they define a relationship between two sets). Like sets, implementations may vary greatly, with the potential to increase efficiency of relation-specific operations.

% Array of tuples, bimaps, bimaps with lists, bimaps with pointers, pathmaps, tries?

Like sets, relations may benefit from compile-time computed properties: size of the domain, size of the range, totality, many-to-oneness, $O(1)$ access to the domain and range, and access to associated domain/range properties.


% Once everything is settled, see about theorem provers, HOL/Lean, to force correct handling of definedness


% TODO add a section discussion nondeterminism, operator behaviour, justification, etc.

\paragraph{Relational Subtypes}
The notion of a relation is a broad classification of a set of pairs, where the domain, range, and relationship are not bound by any conditions. However, when generating code based on purely mathematical relations, it can be advantageous to use notions of totality, injectivity, surjectivity, etc. In particular, these named relational subtypes can be described with four properties: total on the domain, total on the range, one-to-many, and many-to-one.

The below table converts conventional notation into these four subtypes:
\begin{table}[H]
  \centering\begin{tabular}{|c|c|c|c|c|}
    \hline
    Relational Subtype & \multicolumn{4}{c|}{Properties}\\
    & T & TR & 1- & 1-R\\
    \hline
    Relation & & & & \\
     & & & & \checkmark \\
    Partial Function & & & \checkmark & \\
    Partial Injection & & & \checkmark & \checkmark \\
    Surjective Relation & & \checkmark & & \\
     & & \checkmark & & \checkmark \\
    Partial Surjection & & \checkmark & \checkmark & \\
     & & \checkmark & \checkmark & \checkmark \\
    Total Relation & \checkmark & & & \\
     & \checkmark & & & \checkmark \\
    Total Function & \checkmark & & \checkmark & \\
    Total Injection & \checkmark & & \checkmark & \checkmark \\
    Total Surjective Relation & \checkmark & \checkmark & & \\
     & \checkmark & \checkmark & & \checkmark \\
    Total Surjection & \checkmark & \checkmark & \checkmark & \\
    Bijection & \checkmark & \checkmark & \checkmark & \checkmark \\
    % Relation & \\
    % Total Relation & T \\
    % Surjective Relation & TR \\
    % Total Surjective Relation & T, TR \\
    % Partial Function & 1-\\
    % Total Function & T, 1-\\
    % Partial Injection & 1-, 1-R\\
    % Total Injection & T, 1-, 1-R\\
    % Partial Surjection & TR, 1-\\
    % Total Surjection & T, TR, 1-\\
    % Bijection & T, TR, 1-, 1-R\\
    % Should we consider Partial Bijection too?
    % Partial Bijection  & \\
    \hline
  \end{tabular}
  \caption{Conventional mathematical notation of relational subtypes decomposed into four properties: domain totality (T), range totality (TR), domain one-to-manyness (1-), and range one-to-manyness (1-R)}
  \label{tab:relSubtype}
\end{table}

\paragraph{Subtype Properties}
Properties that can affect code generation may involve:
\begin{itemize}
  \item New properties after applying operations on sets/relations
  \item Size of result
  \item Validity of operations
  \item Deterministic nature of operations
  \item Super/subset properties
  \item Relational identities
\end{itemize}
%
Below is a list of the impact of these properties.
\\
Domain Totality:
\begin{itemize}
  \item $R[S] \neq \emptyset$
  \item $R(S)$ is always valid
  \item $R^{-1}$ is TR
  \item $|R| \geq dom(R)$
\end{itemize}
%
Range Totality:
\begin{itemize}
  \item $R^{-1}$ is T
  \item $|R| \geq ran(R)$
\end{itemize}
%
Domain One-to-manyness:
\begin{itemize}
  \item $R(\{e\})$ is deterministic
  \item $|R[\{e\}]| = 1$
  \item $|R[S]| \leq |S|$
  \item $|R| \leq dom(R)$
\end{itemize}
%
Range One-to-manyness:
\begin{itemize}
  \item $|R| \leq ran(R)$
\end{itemize}

\subsubsection{Bags}
Bags (or Multisets) are a refinement of relations wherein, for a bag with element type $T$, the bag is a relational function from $T \rightarrow \mathbb{Z}^+$, denoting the count of items in a collection. Once the count of an item reaches 0, the item is removed from the bag.

Bags must be many-to-one and carry specialized behaviours for conventional relation operations. For example, relational image ($R[B]$) with a bag argument will now return a bag, where the number of occurrences of the resulting item replaces what would otherwise be a regular set of those items.

Additional fields to record the size of the bag (sum of the range) may prove useful.

\subsubsection{Sequences}
Sequences are yet another refinement of relations where a sequence with element type $T$ represents a relational function from $\mathbb{N} \rightarrow T$ that is total on some defined bounds.

The dense nature of sequence domains, along with knowledge of sequence sizes, may open up further data type refinements.

\subsection{Data Type Operations}
\subsubsection{Primitives}
\subsubsection{Sets}
\subsubsection{Relations}
\subsubsection{Bags}
\subsubsection{Sequences}

\section{Implementation Specification}
\subsection{Abstract Data Types}
\subsection{Atomic Operations}
\subsection{Optimization Frontend}
\subsection{Optimization Backends}

\section{AST Lowering and Optimization Rules}
Below is a list of rewrite rules for key abstract data types, syntactic sugar, and some builtin functions. Phases are intended to be executed in order; the post-condition of one phase serves as the pre-condition for the next. Standard functions (like map, flatten, length, etc.) are written as a library in Simile rather than baked-in optimizations.


\subsection{Syntactic Sugar for Bags}
\begin{align}
  \tag{Bag Image}
  R[B]
  &\leadsto
  R^{-1} \circ B
  \\
  \tag{Bag Predicates - Union}
  S \cup T &\leadsto \bag{x \in dom(S) \cup dom(T)}{n = max(S[x] \cup T[x])}{x \mapsto n}
  \\
  \tag{Bag Predicates - Intersection}
  S \cap T &\leadsto \bag{x \in dom(S)}{x \in dom(T) \land n = min(S[x] \cup T[x]) \land n > 0}{x \mapsto n}
  \\
  \tag{Bag Predicates - Sum}
  S + T &\leadsto \bag{x \in dom(S) \cup dom(T)}{n = sum(S[x] \cup T[x])}{x \mapsto n}
  \\
  \tag{Bag Predicates - Difference}
  % S - T &\leadsto \bag{x}{x \in dom(S) \land pop(S[x]) - popDefault(T[x], 0) \land n > 0}{x \mapsto n}
  S - T &\leadsto \bag{x \in dom(S)}{n = S(x) - (T \cup \llbracket x \mapsto 0\rrbracket)(x) \land n > 0}{x \mapsto n}
\end{align}

\subsection{Syntactic Sugar for Sequences}
\begin{align}
    \tag{Concatenation}
    L \concat M
    &\leadsto
    L \cup \Set{x \mapsto y \in M }{x + max(dom(L)) \mapsto y}
\end{align}

\subsection{Syntactic Sugar for Relations}
\noindent\begin{minipage}{\linewidth}
\begin{align}
  \tag{Functional Override}
  f(x) := E
  &\leadsto
  f := f \oplus \{x \mapsto E\}
  \\
  \tag{Override}
  R \oplus Q &\leadsto Q \cup (dom(Q) \domsub R)
  \\
  \tag{Domain Restriction}
  S \triangleleft R &\leadsto \bSet{x \mapsto y \in R}{x \in S}{x \mapsto y}
  \\
  \tag{Domain Subtraction}
  S \domsub R &\leadsto \bSet{x \mapsto y \in R}{x \notin S}{x \mapsto y}
  \\
  \tag{Range Restriction}
  R \triangleright S &\leadsto \bSet{x \mapsto y \in R}{y \in S}{x \mapsto y}
  \\
  \tag{Range Subtraction}
  R \ransub S &\leadsto \bSet{x \mapsto y \in R}{y \notin S}{x \mapsto y}
\end{align}
\end{minipage}

\subsection{Comprehension Construction}

\paragraph{Intuition} All set-like variables and literals are decomposed into set comprehensions.

\paragraph{Post-condition}
\begin{itemize}
  \item All terms with a set-like type (relations, bags, sets, sequences, etc.) must be in comprehension form.
\end{itemize}


\noindent\begin{minipage}{\linewidth}
\begin{align}
  \tag{Functional Image \footnote{May fail if R is not total. $choice$ is nondeterministic.}}
  R(x) &\leadsto choice(R[x])
  \\ % dont use pop - random choice but pop(S) == pop(S). Python op is completely nondeterministic so it may not return the same one R(x) = small epsilon y . y in R[{x}]
  % if R[{x}] is empty, you return "an arbitrary element of range of R", not None
  \tag{Image}
  R[S] &\leadsto \bSet{x \mapsto y \in R}{x \in S}{y}
  \\
  \tag{Product}
  x \mapsto y \in S \times T &\leadsto x \in S \land y \in T
  \\
  \tag{Inverse}
  x \mapsto y \in R^{-1} &\leadsto y \mapsto x \in R
  \\
  \tag{Composition}
  x \mapsto y \in (Q \circ R) &\leadsto x \mapsto z \in Q \land z' \mapsto y \in R \land z=z'
  % \tag{Exists Elimination}
  % \exists z \cdot x \mapsto z \in Q \land z \mapsto y \in R &\leadsto x \mapsto z \in Q \land z' \mapsto y \in R \land z=z'
  % \\
  \end{align}
\end{minipage}
\noindent\begin{minipage}{\linewidth}
\begin{align}
  \tag{Predicate Operations - Union}
  S \cup T
  &\leadsto
  \setOne{(x \in S \mid x) ` (x \in T \cdot x \notin S \mid x)}
  \\
  \tag{Predicate Operations - Intersection}
  S \cap T
  &\leadsto
  \Set{x \in S}{x \in T}
  \\
  \tag{Predicate Operations - Difference}
  S \setminus T
  &\leadsto
  \Set{x \in S}{x \notin T}
  % \\
  % \tag{Singleton Membership \footnote{Currently unused. We need to be careful to handle the case where $x$ is a free variable.}}
  % x \in \{e\}
  % &\leadsto
  % x = e
  \\
  \tag{Membership Collapse \footnote{Rule only matches inside the predicate of a quantifier. Explicitly enumerating all matches for all quantuantification types and predicate cases (ANDs, ORs, etc.) would require too much boilerplate. $x$ must be bound by the encasing quantifier.\\
  The $\oplus$ operator represents any quantifier that returns a set-like type (ex. generalized union/intersection, set comprehension, relation comprehension).}}
  x \in \oplus(E \mid P)
  &\leadsto
  P \land x = E
  \\
  \tag{Set Membership}
  e \in S
  &\leadsto
  \exists{x \in S \mid e = x}
\end{align}
\end{minipage}

\subsection{Disjunctive Normal Form}

\paragraph{Intuition} All quantifier predicates are expanded to DNF (i.e. $\land$-operations nested within top-level $\lor$-operations).

\paragraph{Notes} All matches of this phase occur only inside quantifier predicates.

\paragraph{Post-condition}
\begin{itemize}
  \item All terms with a set-like type (relations, bags, sets, sequences, etc.) must be in comprehension form.
  \item Quantifier predicates are in disjunctive normal form - top level or-clauses with inner and-clauses.
  \item If no $\lor$ operators exist within a quantifier predicate, the predicate must only contain $\land$ operators.
\end{itemize}

% https://en.wikipedia.org/wiki/Disjunctive_normal_form
\noindent\begin{minipage}{\linewidth} % Need minipage for footnote
\begin{align}
  \tag{Flatten Nested $\land$}
  x_1 \land ... \land (x_i \land x_{i+1}) \land ...
  &\leadsto
  x_1 \land ... \land x_i \land x_{i+1} \land ...
  \\
  \tag{Flatten Nested $\lor$}
  x_1 \lor ... \lor (x_i \lor x_{i+1}) \lor ...
  &\leadsto
  x_1 \lor ... \lor x_i \lor x_{i+1} \lor ...
  \\
  \tag{Double Negation}
  \lnot \lnot x
  &\leadsto
  x
  \\
  \tag{Distribute De Morgan - Or}
  \lnot (x \lor y)
  &\leadsto
  \lnot x \land \lnot y
  \\
  \tag{Distribute De Morgan - And}
  \lnot (x \land y)
  &\leadsto
  \lnot x \lor \lnot y
  \\
  \tag{Distribute $\land$ over $\lor$}
  x \land (y \lor z)
  &\leadsto
  (x \land y) \lor (x \land z)
\end{align}
\end{minipage}

\subsection{Or-wrapping}
\paragraph{Intuition} Restructuring quantifier predicates with top-level ORs for later rules.

\paragraph{Post-condition}
\begin{itemize}
  \item All terms with a set-like type (relations, bags, sets, sequences, etc.) must be in comprehension form.
  \item Quantifier predicates are in disjunctive normal form - top level or-clauses with inner and-clauses.
\end{itemize}

\noindent\begin{minipage}{\linewidth}
\begin{align}
  \tag{Or-wrapping \footnote{To simplify the matching process later on, we wrap every top-level AND statement (which is guaranteed to be a ListOp by the dataclass field type definition) with an OR.}}
  \bSetT{E}{\bigwedge P_i}
  &\leadsto
  \bSetT{E}{\bigvee \bigwedge P_i}
\end{align}
\end{minipage}

\subsection{Generator Selection}
\paragraph{Intuition} All nested ANDs are placed in a tree structure to better suit loop lowering.

\paragraph{Post-condition}
\begin{itemize}
  \item All terms with a set-like type (relations, bags, sets, sequences, etc.) must be in comprehension form.
  \item Each top-level or-clause within quantification predicates must have one selected `generator' predicate (GeneratorSelectionPredicate - GSP) of the form $x \in S$ that loops over its bound dummy variable $x$.
  \item All conjunctive clauses are wrapped in either a GSP or CombGSP structure.
  \item Quantifier predicates are in disjunctive normal form (with GSP replacing AND structures) - top level or-clauses with inner and-clauses.
\end{itemize}

\paragraph{Notes} All matches of this phase occur only inside quantifier predicates. GSP is a tuple of form (generator chain, predicates) that can be flattened to a conjunctive clause. A CombGSP is a triple of form (shared generator, disjunctive children predicates/generators, predicates) that can likewise be flattened into a conjunctive clause. We keep the structure of these tuples distinct to ease the rewriting process.

\paragraph{Brainstorming selection heuristics}
Conditions to consider:
\begin{itemize}
  \item Prefer composition chains in case of nested loops (ex. $x \mapsto y \in R \land y \mapsto z \in Q$, but what about renaming? - $x \mapsto y \in R \land y' \mapsto z \in Q \land y = y'$). Perhaps this could just be a preference for relations over sets if we see a nested quantifier.
  \item Choose most common generator among a list of or-clauses (allows for greater simplification later on)
  \item Choose smallest generator (by set size). This information may not always be accessible (ex. if a set is created by reading values from standard input). Functions may need to choose this on a case-by-case basis (ie. one function call could have two args, with a small set on the left arg and a large set on the right arg. But what if the sizes are reversed later in the code? We've already statically lowered it).
  \item When should constant folding happen? Equality substitution necessary for this?
  \item This may need to work with nesting considerations.
  \item What if we try moving these optimizations to $loop$ structures far later in the pipeline?
\end{itemize}

\noindent\begin{minipage}{\linewidth}
\begin{align}
  \tag{GSP Wrapping \footnote{GSP is a tuple-like structure, where all entries can be flattened into an AND. The LH term must occur inside a quantifier's predicate - one generator per top-level or-clause. Chained generators may be necessary to allow for nested loop operations (ie. non-top-level clauses may require generators). $P_g$ is a list of chained generators, specific set membership clauses (of form $x \in S$) distinguished from the rest of $\bigwedge P_i$. Currently, the selection of $P_g$ is arbitrary (and thus the rewrite system is not confluent), but heurestics may be added later to choose better generators.}}
  \bigwedge P_i
  &\leadsto
  GSP(\bigwedge P_{g_i}, \bigwedge_{P_i \notin P_g} P_i)
  \\
  \tag{Nested Generator Selection \footnote{$free$ is a function that returns true when variables in the predicate are defined (either bound by the current generator or bound prior to the current statement)}}
  GSP(x \land xs, P) \lor GSP(x \land ys, Q)
  &\leadsto
  CombGSP(x, GSP(xs, P) \lor GSP(ys, Q))
\end{align}
\end{minipage}



\subsection{GSP to Loop}

\paragraph{Intuition} Start lowering expressions into imperative-like loops.

\paragraph{Post-condition}
\begin{itemize}
  \item All quantifiers are transformed into $loop$ structures.
  \item $loop$ predicates are of the form $x \in S \land \bigwedge P_i$.
  \item All $loop$ predicates have an assigned generator of the form $x \in S$.
\end{itemize}

\noindent\begin{minipage}{\linewidth}
\begin{align}
  \tag{Quantifier Generation \footnote{$\oplus$ works for any quantifier (but not $\forall$ and $\exists$). The identity and accumulate functions are determined by the realized $\oplus$. For example, if $\oplus = \sum$, the identity is 0 and accumulate is addition.}}
  \oplus E \mid P
  &\leadsto
  \begin{minipage}[]{0.4\textwidth}
  \begin{algorithmic}
  \State $a := identity(\oplus)$
    \Loop\ $P:$ % should we use the \in operator here, or just plaintext in?
        \State $a := accumulate(a, E)$
    \EndLoop
  \end{algorithmic}
  \end{minipage}
\end{align}
\end{minipage}
% \noindent\begin{minipage}{\linewidth}
% \begin{align}
%   \tag{Nested Quantifier}
%   \begin{minipage}[]{0.5\textwidth}
%   \begin{algorithmic}
%     \State $a := accumulator(a, \oplus(E \mid P))$
%   \end{algorithmic}
%   \end{minipage}
%   &\leadsto
%   \begin{minipage}[]{0.4\textwidth}
%   \begin{algorithmic}
%     \Loop\ $P:$
%       \State $a := accumulator(a, E)$
%     \EndLoop
%   \end{algorithmic}
%   \end{minipage}
% \end{align}
% \end{minipage}
\noindent\begin{minipage}{\linewidth}
\begin{align}
  \tag{Top-level Or-Loop}
  \begin{minipage}[]{0.3\textwidth}
  \begin{algorithmic}
    \Loop\ $\bigvee P_i:$
      \State body
    \EndLoop
  \end{algorithmic}
  \end{minipage}
  &\leadsto
  \begin{minipage}[]{0.5\textwidth}
  \begin{algorithmic}
    \Loop\ $P_0$
      \State $body$
    \EndLoop
    \Loop\ $P_1 \land \lnot P_0$
      \State $body$
    \EndLoop
    \Loop\ $P_2 \land \lnot \bigvee_{i < 2} P_i$
      \State $body$
    \EndLoop
  \end{algorithmic}
  \end{minipage}
\end{align}
\end{minipage}
%
\noindent\begin{minipage}{\linewidth}
\begin{align}
  \tag{Chained GSP Loop \footnote{$free$ considers defined variables at this point in time.}}
  \begin{minipage}[]{0.4\textwidth}
  \begin{algorithmic}
    \Loop\ $GSP(g \land gs, P)$
      \State body
    \EndLoop
  \end{algorithmic}
  \end{minipage}
  &\leadsto
  \begin{minipage}[]{0.4\textwidth}
  \begin{algorithmic}
    \Loop\ $SingleGSP(g, free(P))$
      \Loop\ $GSP(gs, bound(P))$
        \State body
      \EndLoop
    \EndLoop
  \end{algorithmic}
  \end{minipage}
\end{align}
\end{minipage}
%
% \noindent\begin{minipage}{\linewidth}
% \begin{align}
%   \tag{Single GSP Loop}
%   \begin{minipage}[]{0.4\textwidth}
%   \begin{algorithmic}
%     \Loop\ $GSP(g, P)$
%       \State body
%     \EndLoop
%   \end{algorithmic}
%   \end{minipage}
%   &\leadsto
%   \begin{minipage}[]{0.4\textwidth}
%   \begin{algorithmic}
%     \Loop\ $SingleGSP(g, P)$
%       \State body
%     \EndLoop
%   \end{algorithmic}
%   \end{minipage}
% \end{align}
% \end{minipage}
%
\noindent\begin{minipage}{\linewidth}
\begin{align}
  \tag{Empty GSP Loop}
  \begin{minipage}[]{0.4\textwidth}
  \begin{algorithmic}
    \Loop\ $GSP([], P)$
      \State body
    \EndLoop
  \end{algorithmic}
  \end{minipage}
  &\leadsto
  \begin{minipage}[]{0.4\textwidth}
  \begin{algorithmic}
    \If{$P$}
      \State body
    \EndIf
  \end{algorithmic}
  \end{minipage}
\end{align}
\end{minipage}
%
\noindent\begin{minipage}{\linewidth}
\begin{align}
  \tag{Combined GSP Loop}
  \begin{minipage}[]{0.4\textwidth}
  \begin{algorithmic}
    \Loop\ $CombGSP(g, gs, P)$
      \State body
    \EndLoop
  \end{algorithmic}
  \end{minipage}
  &\leadsto
  \begin{minipage}[]{0.4\textwidth}
  \begin{algorithmic}
    \Loop\ $SingleGSP(g, P)$
      \Loop\ $gs_0$
        \State body
      \EndLoop
      \Loop\ $gs_1$
        \State body
      \EndLoop
      \State ...
    \EndLoop
  \end{algorithmic}
  \end{minipage}
\end{align}
\end{minipage}

% \noindent\begin{minipage}{\linewidth}
% \begin{align}
%   \tag{Disjunct conditional}
%   \begin{minipage}[]{0.33\textwidth}
%   \begin{algorithmic}
%     \Loop\ $\bigvee GSP_i(P_{g_i}, P_i)$
%       \State body
%     \EndLoop
%   \end{algorithmic}
%   \end{minipage}
%   &\leadsto
%   \begin{minipage}[]{0.5\textwidth}
%   \begin{algorithmic}
%     \Loop\ $GSP_0(P_{g_0}, P_0):$
%       \State body
%     \EndLoop
%     \Loop\ $GSP_1(P_{g_1}, P_1 \land \lnot GSP_0):$
%       \State body
%     \EndLoop
%     \Loop\ $GSP_2(P_{g_2}, P_2 \land \lnot (\bigvee_{i < 2} GSP_i)):$
%       \State body
%     \EndLoop
%     \State ...
%   \end{algorithmic}
%   \end{minipage}
%   \\
%   \tag{Disjunct Conditional 2}
%   \begin{minipage}[]{0.3\textwidth}
%   \begin{algorithmic}
%     \Loop\ $GSP(P_g, \bigvee P_i)$
%       \State body
%     \EndLoop
%   \end{algorithmic}
%   \end{minipage}
%   &\leadsto
%   \begin{minipage}[]{0.8\textwidth}
%   \begin{algorithmic}
%     \Loop\ $GSP(P_g, \bigwedge True):$
%       \If{$\bigvee_{ P_i \neq GSP(...)} P_i$}
%         \State body
%       \EndIf
%       \Loop\ $\bigvee_{P_i: GSP(P_{ig}, P_{is})} GSP(P_{g_i}, P_{is} \land \lnot \bigvee_{P_i \neq GSP(...)} P_i):$
%         \State body
%       \EndLoop
%     \EndLoop
%   \end{algorithmic}
%   \end{minipage}
% \end{align}
% \end{minipage}

\subsection{Relational Subtyping Loop Simplification}

\paragraph{Intuition} Simplify unnecessary iteration structures into direct accesses.

\paragraph{Post-condition}
\begin{itemize}
  \item All quantifiers are transformed into $loop$ structures.
  \item $loop$ predicates are of the form $x \in S \land \bigwedge P_i$.
  \item All $loop$ predicates have an assigned generator of the form $x \in S$.
\end{itemize}
No change from previous, but the number of loop structures should not increase.

\noindent\begin{minipage}{\linewidth}
\begin{align}
  \tag{Total membership elimination \footnote{$R$ is total and $x$ satisfies the type of $R$'s domain. Currently unused since it may eliminate some generators - move higher in the list?}}
  x \in dom(R)
  &\leadsto
  True
  \\
  \tag{Surjective membership elimination \footnote{$R$ is surjective and $x$ satisfies the type of $R$'s range. Currently unused since it may eliminate some generators - move higher in the list? Or }}
  x \in ran(R)
  &\leadsto
  True
\end{align}
\end{minipage}
\noindent\begin{minipage}{\linewidth}
\begin{align}
  \tag{Concrete Domain Image}
  \begin{minipage}[]{0.55\textwidth}
  \begin{algorithmic}
    \Loop\ $SingleGSP(x \mapsto y \in R, x = a \land P)$
      \State body
    \EndLoop
  \end{algorithmic}
  \end{minipage}
  &\leadsto
  \begin{minipage}[]{0.6\textwidth}
  \begin{algorithmic}
    \If{$a \in dom(R)$}
      \Loop\ $SingleGSP(y \in R[a], P[x := a])$
        \State body
      \EndLoop
    \EndIf
  \end{algorithmic}
  \end{minipage}
\end{align}
\end{minipage}
\noindent\begin{minipage}{\linewidth}
\begin{align}
  \tag{Concrete Range Image \footnote{Requires efficient lookups for $ran$, inverse, and image}}
  \begin{minipage}[]{0.55\textwidth}
  \begin{algorithmic}
    \Loop\ $SingleGSP(x \mapsto y \in R, y = a \land P)$
      \State body
    \EndLoop
  \end{algorithmic}
  \end{minipage}
  &\leadsto
  \begin{minipage}[]{0.6\textwidth}
  \begin{algorithmic}
    \If{$a \in ran(R)$}
      \Loop\ $SingleGSP(x \in R^{-1}[a], P[y := a])$
        \State body
      \EndLoop
    \EndIf
  \end{algorithmic}
  \end{minipage}
\end{align}
\end{minipage}
\noindent\begin{minipage}{\linewidth}
\begin{align}
  \tag{Single Element Loop}
  \begin{minipage}[]{0.5\textwidth}
  \begin{algorithmic}
    \Loop\ $SingleGSP(y \in R[x], P)$
      \State body
    \EndLoop
  \end{algorithmic}
  \text{where $free(x)$ and $R$ is many-to-one}
  \end{minipage}
  &\leadsto
  \begin{minipage}[]{0.5\textwidth}
  \begin{algorithmic}
    \If{$P[y:=R(a)]$}
      \State $body[y:=R(a)]$
    \EndIf
  \end{algorithmic}
  \end{minipage}
\end{align}
\end{minipage}
\noindent\begin{minipage}{\linewidth}
\begin{align}
  \tag{Singleton Membership Elimination}
  \begin{minipage}[]{0.5\textwidth}
  \begin{algorithmic}
    \Loop\ $SingleGSP(x \in \{y\}, P)$
      \State body
    \EndLoop
  \end{algorithmic}
  \end{minipage}
  &\leadsto
  \begin{minipage}[]{0.5\textwidth}
  \begin{algorithmic}
    \If{$P[x:=y]$}
      \State $body[x:=y]$
    \EndIf
  \end{algorithmic}
  \end{minipage}
\end{align}
\end{minipage}

\subsection{Loop Code Generation}

\paragraph{Intuition} Eliminate all intermediate loop structures.

\paragraph{Post-condition}
\begin{itemize}
  \item AST is in imperative code style (for loops, if statements, etc).
  \item All $loop$ and quantification constructs have been eliminated.
  \item Some variables may not be defined.
\end{itemize}

\noindent\begin{minipage}{\linewidth}
\begin{align}
  \tag{Conjunct Conditional \footnote{Function $free$ returns clauses in $P$ that contain only free + defined variables. $bound$ returns the clauses that contain the bound variable $x$ or undefined variables. $P_g$ is the selected generator.}}
  \begin{minipage}[]{0.4\textwidth}
  \begin{algorithmic}
    \Loop\ $SingleGSP(P_g, \bigwedge P_i)$
      \State body
    \EndLoop
  \end{algorithmic}
  \end{minipage}
  &\leadsto
  \begin{minipage}[]{0.5\textwidth}
  \begin{algorithmic}
    \If{$\bigwedge_{free(P_i)} P_i$}
      \For{$P_g$}
        % \State $P_as$ % TODO rewrite these to be more clear that they are assignments
        \If{$\bigwedge_{bound(P_i)} P_i$}
            \State body
          % \If{$Q_0$}
          %   \State body
          % \EndIf
          % \If{$Q_1 \land \lnot Q_0$}
          %   \State body
          % \EndIf
          % \If{$Q_2 \land \lnot (\bigvee_{i < 2} Q_i)$}
          %   \State body
          % \EndIf
          % \State ...
        \EndIf
      \EndFor
    \EndIf
  \end{algorithmic}
  \end{minipage}
\end{align}
\end{minipage}



\subsection{Replace and Simplify}

\paragraph{Intuition} Eliminate all undefined variables (with implicit $\exists$ quantifiers).

\paragraph{Post-condition}
\begin{itemize}
  \item AST is in imperative code style (for loops, if statements, etc).
  \item All variables are defined.
\end{itemize}

\noindent\begin{minipage}{\linewidth}
\begin{align}
  \tag{Equality Elimination \footnote{$x$ is an undefined, unbound variable in the current scope. $E$ is an expression that does not contain $x$. Simplify resulting booleans.}}
  \begin{minipage}[]{0.5\textwidth}
  \begin{algorithmic}
    \If{$Identifier(x) = E \land P$}
      \State body
    \EndIf
  \end{algorithmic}
  \end{minipage}
  &\leadsto
  \begin{minipage}[]{0.4\textwidth}
  \begin{algorithmic}
    \If{$P[x := E]$}
      \State body[x := E]
    \EndIf
  \end{algorithmic}
  \end{minipage}
\end{align}
\end{minipage}
%
\noindent\begin{minipage}{\linewidth}
\begin{align}
  \tag{Simplify And-true}
  P \land True
  &\leadsto
  P
  \\
  \tag{Simplify And-false}
  P \land False
  &\leadsto
  False
  \\
  \tag{Simplify Or-true}
  P \lor True
  &\leadsto
  True
  \\
  \tag{Simplify Or-false}
  P \lor False
  &\leadsto
  P
  \\
  \tag{Flatten Nested Statements}
  Statement(Statement(x))
  &\leadsto
  Statement(x)
\end{align}
\end{minipage}


% \section{Implementation Representation}
% Different implementations of each data type will have varying strengths and weaknesses, not only in theoretical asymptotic time and space, but in concrete real-world tests. Cache usage and additional information through object metadata may prove influential on smaller tests. Since this document is only concerned with the theoretical compiler specification, we analyze the theoretical time and space complexity, then pair gathered examples with a test plan for hardware considerations.

% A first approach to tackling these type representations would likely constitute a linked list. The space requirements for enumeration are straightforward, with extra allocations for link pointers. Insertions for unordered collections or append/concat operations are $O(1)$, but $O(n)$ for indexed insertion and union with one element. Lookups for all collections are $O(n)$, but this running time is undesirable for the often-used \texttt{in} operator for set-generated code. Since linked lists naturally enforce element order, this structure may be suitable for fast-changing sequences. Although, a limited-size sequence may be better suited for a contiguous array for $O(1)$ indexing. \textit{TODO: For sequences, we should also see if trees/heaps or bloom filters could provide efficient membership checking. Bloom filters are probabilistic but can determine $\neq$ operations.}

% On the other hand, hashmaps with $O(1)$ membership and element lookups are useful for all unordered collections. Relations may need bidirectional hashmaps that can efficiently handle many-to-many relations.

% Compressed bitmaps may be used for sets, but require a lot of space for sparse elements.

% Bags may be implemented either as a (linked) list, a set of tuples where the number of element occurrences is stored in the second tuple component, or a relation where the number of occurrences is the codomain.
\printbibliography

\end{document}